% !TEX program = xelatex

%\documentclass[12pt, letterpaper]{article}

\documentclass[14pt]{extarticle}   % also supports 8,9,10,12,14,17,20pt
\usepackage[a4paper,margin=0.5in]{geometry} % adjust margin as needed


\usepackage{fontspec}
\usepackage{polyglossia}
\usepackage{setspace}

% Main and other language settings
\setmainlanguage{sanskrit}
\setotherlanguage{english}

% Fonts
\newfontfamily\sanskritfont[Script=Devanagari, 
  % Mapping=devanagarinumerals,
  UprightFont = *-Regular,
  Ligatures=TeX,
  % BoldFont = *-Medium % or *-Medium if available
]{Shobhika}
\newfontfamily\englishfont{Times New Roman} % or any other font you like

\newcommand{\eng}[1]{\begin{english}#1\end{english}}
\newcommand{\emd}{— }

% Tie the sutracount to the section counter
\newcounter{sutracount}[subsection] 


% Update your command to take only 1 argument (the text)
\newcommand{\sutra}[1]{%
  \stepcounter{sutracount}% Increments the counter by 1
  \textbf{#1} ॥\thesutracount॥% Prints the current counter value
}

% Redefine the counter output to include the section number
\renewcommand{\thesutracount}{\thesubsection.\arabic{sutracount}}


\setlength{\parskip}{1em}      % Vertical space between paragraphs
\onehalfspacing    % 1.5x line spacing
% \setlength{\skip\footins}{20pt}  % default is about 10pt, increase as needed


\setlength{\parindent}{0in}

\title{कातन्त्रलघुवृत्तिः}
\author{छुच्छुकभट्टः}
\date{\eng{\today}}

\begin{document}

\maketitle

\begin{center}
  \textbf{प्रकाशकः} --- लोकेश शर्मा
\end{center}

\hrulefill

\begin{center}
    ॐ श्रीगणेशाय नमः॥
\end{center}

\section{सन्धिकरणम्}

\subsection{संज्ञापादः}

\sutra{सिद्धो वर्णसमाम्नायः}

वर्णानाम् अकारादीनां वर्णानां यः समाम्नायः क्रमेण पाठः सः सिद्धः प्रसिद्धः लोके इति किं तेनोपदिष्टेन आविद्वदङ्गनागोपालबालपर्यन्तं यावत्। \\
वर्णाः ४९–\\
अआ इई उऊ ऋॠ ऌॡ एऐ ओऔ अं अः \\
कखगघङ चछजझञ टठडढण तथदधन पफबभम यरलव शषसह। \\
वर्णसंज्ञया क्व प्रयोजनम्? सर्वत्र।

\sutra{तत्र चतुर्दशादौ स्वराः}

तत्र प्रसिद्धे वर्णसमाम्नाये आदौ ये चतुर्दश वर्णाः ते स्वरसंज्ञा भवन्ति। \\
स्वराः १४–\\
अआ इई उऊ ऋॠ ऌॡ एऐ ओऔ\\
स्वरसंज्ञया क्व प्रयोजनम्? रो रे लोपं स्वरश्च पूर्वो दीर्घ इत्यादिषु।

\sutra{दश समानाः}

तत्र प्रसिद्धे वर्णसमाम्नाये आदौ ये दशवर्णाः ते समानसंज्ञा भवन्ति। \\
समानाः १०–\\
अआ इई उऊ ऋॠ ऌॡ \\
समानसंज्ञया क्व प्रयोजनम्? समानः सवर्णे दीर्घीभवति परश्चलोपमित्यादिषु।

\sutra{तेषां द्वौ द्वावन्योऽन्यस्य सवर्णौ}
तेषां दशानां समानानां मध्यात् यौ द्वौ द्वौ वर्णौ तौ अन्योऽन्यस्य परस्परं सवर्णसंज्ञौ भवतः। \\
सवर्णाः १०–\\
अआ इई उऊ ऋॠ ऌॡ\\
कः कस्य सवर्णः? अकास्य आकारः सवर्णः। आकारस्य अकारः सवर्णः। उकारस्य ऊकारः सवर्णः। ऊकारस्य उकारः सवर्णः। ऋकारस्य ॠकारः सवर्णः। ॠकारस्य ऋकारः सवर्णः। ऌकारस्य ॡकारः सवर्णः। ॡकारस्य ऌकारः सवर्णः। \\
सवर्णसंज्ञया क्व प्रयोजनम्? समानः सवर्णे दीर्घीभवति परश्चलोपमित्यादिषु।

\sutra{पूर्वो ह्रस्वः}

तयोः द्वयोर् द्वयोः सवर्णयोः यो यः पूर्वो वर्णः स ह्रस्वसंज्ञो भवति। \\
ह्रस्वाः ५ – अइउऋऌ। \\
ह्रस्वसंज्ञया क्व प्रयोजनम्? ह्रस्वो स्वार्थानामित्यादिषु।

\sutra{पूर्वस्य लघुसंज्ञापि वक्तव्या}

लघवः ५ – अइउऋऌ। \\
लघुसंज्ञया क्व प्रयोजनम्? नामिनश्चापधायालघोरित्यादिषु।

\sutra{परो दीर्घः}

तयोः द्वयोर् द्वयोः सवर्णयोः यो यः परः स दीर्घसंज्ञो भवति। \\
दीर्घाः ५ – आईऊॠॡ। \\
दीर्घसंज्ञया क्व प्रयोजनम्? दीर्घमामिसनावित्यादिषु।

\sutra{परस्य गुरुसंज्ञापि वक्तव्या}

गुरवः ५ - आईऊॠॡ। \\
गुरुसंज्ञया क्व प्रयोजनं गरोश्चनिष्टायां सेट इत्यादिषु।

\sutra{स्वरो ऽवर्णवर्जो नामी}

अवर्णः अकाराकारौ तौ वर्जयित्वा शिष्टः स्वरः नामिसंज्ञो भवति। \\
नामिनः १२ – इई उऊ ऋॠ ऌॡ एऐ ओऔ। \\
नामिसंज्ञया क्व प्रयोजनम्? नामिनः स्वरे इत्यादिषु।

\sutra{एकारादीनि सन्ध्यक्षराणि}

एकारादीनि वर्णानि सन्ध्यक्षरसंज्ञा भवन्ति। \\
सन्ध्यक्षराणि ४ – एऐओऔ। \\
सन्ध्यक्षसंज्ञया क्व प्रयोजनम्? सन्ध्यक्षरान्तानाम् आकारो विकरण इत्यादिषु।

\sutra{कादीनि व्यञ्जनानि}

ककारादीनि वर्णानि व्यञ्जनसंज्ञानि भवन्ति। \\
व्यञ्जनानि ३३ – कखगघङ चछजझञ टठडढण तथदधन पफबभम यरलव शषसह। \\
व्यञ्जनसंज्ञया क्व प्रयोजनं मोऽनुस्वारं व्यञ्जन इत्यादिषु।

\sutra{ते वर्गाः पञ्चपञ्चशः}

\textbf{ते} कादयो वर्णाः \textbf{पञ्चपञ्चशः} पञ्चपञ्चाक्षराणि भूत्वा \textbf{वर्गा} उच्यन्ते। 
वर्गाः ५ – कचटतप। \\
वर्ग्याः २५ – \\
कवर्गः - कखगघङ\\
चवर्गः - चछजझञ\\
टवर्गः - टठडढण\\
तवर्गः - तथदधन\\
पवर्गः - पफबभम।\\
वर्गसंज्ञया क्व प्रयोजनम्? वर्ग्ये तद्वर्गपञ्चमं वेत्यादिषु।

\sutra{वर्गाणां प्रथमद्वितीयाः शषसाश्चाघोषाः}

वर्गाणां ये प्रथमद्वितीयावर्णाः तथा शषसाश्च ते अघोषसंज्ञा भवन्ति। \\
अघोषाः १३ – कख चछ टठ तथ पफ शषस। \\
अघोषसंज्ञया क्व प्रयोजनम्? अघोषे प्रथम इत्यादिषु।।

\sutra{घोषवन्तोऽन्ये}

अघोषेभ्यः येऽन्ये वर्णास् ते घोषवत्संज्ञा भवन्ति। \\
घोषवन्तः २० – गघङ जझञ डढण दधन बभम यरलवह। \\
घोषवत्संज्ञया क्व प्रयोजनम्? अकारो दीर्घं घोषवतीत्यादिषु।।

\sutra{अनुनासिका ङञणनमाः}

ङञणनमा वर्णाः अनुनासिकसंज्ञा भवन्ति। \\
अनुनासिका ५ – ङञणनम। \\
अनुनासिकसंज्ञया क्व प्रयोजनम्? अतोन्तोनुस्वारोनुनासिकान्तस्येत्यादिषु।।

\sutra{अन्तःस्था यरलवाः}

यरलवा वर्णाः अन्तस्थसंज्ञा भवन्ति। \\
अन्तःस्थाः ४ – यरलव। \\
अन्तःस्थसंज्ञया क्व प्रयोजनम्? स परस्वरायासंप्रसारणमन्तःस्थाया इत्यादिषु।।

\sutra{ऊष्माणः शषसहाः}

शषसह वर्णा ऊष्मसंज्ञा भवन्ति। \\
ऊष्माणः – ४ शषसह। \\
ऊष्मसंज्ञायाः प्रयोजनं नास्तीत्यपार्थकोयं योगारम्भ इत्याहुः।

\sutra{अः इति वसर्जनीयः}

अकारात्परं बिन्दुद्वयं विसर्जनीयसंज्ञं भवति। \\
विसर्जनीयसंज्ञया क्व प्रयोजनम्? विसर्जनीयश्छेछोवाशमित्यादिषु।।

\sutra{ᳳक इति जिह्वामूलीयः}

ककारात् पूर्वोवर्णः जिह्वामूलीयसंज्ञोभवति। जिह्वामूलीयः कीदृशः? यवाकृतिसदृशः। \\
जिह्वामूलीयसंज्ञया क्व प्रयोजनम्? कखयो जिह्वामूलीयं नवेत्यादिषु।।

\sutra{ᳲप इत्युपध्मानीयः}

पकारात्पूर्वो वर्णः उपध्मानीयसंज्ञो भवति। उपध्मानीयः कीदृशः हस्तिकवाटसदृशः। \\
उपध्मानीयसंज्ञया क्व प्रयोजनम्? पफयोरुपध्मानीयं नवेत्यादिषु।।

\sutra{अं इत्यनुस्वारः}

अं इत्यत्रयोबिन्दुः सोsनुस्वारसंज्ञो भवति। \\
अनुस्वारसंज्ञया क्व प्रयोजनम्? मोनुस्वारं व्यञ्जने इत्यादिषु।।

\sutra{पूर्वपरयोर् अर्थोपलब्धौ पदम्}

पूर्वस्य लिङ्गस्य परस्य स्यादेः तथा पूर्वस्य धातोः परस्य त्यादेः यः समुदाय एकीभावः तस्य पदमिति संज्ञा भवति अर्थोपलब्धौ सत्याम्। \\
पदसंज्ञया क्व प्रयोजनं? युष्मदस्मदोः पदं पदात् षष्ठीचतुर्थीद्वितीयासु वस्रसावित्यादिषु।।

\sutra{अनतिक्रमयन् विश्लेषयेत्}

वर्णानाम् एकत्र युक्तानाम् अतिक्रमं विपर्यासं अकुर्वन् विश्लेषयेत्। येनैव क्रमेण वर्णाः संयुज्यन्ते तेनैव क्रमेण वर्णा विश्लिष्यन्ते इति परिभाष्यते। अनियमे नियमः परिभाषा।

\sutra{व्यञ्जनम् अस्वरं परं वर्णं नयेत्}

व्यञ्जनं यत् अस्वरं स्वररहितं तत् परं वर्णं नयेत् प्रापयेत्। त्वमत्र। सत्कर्म। अहमिह।

\sutra{लोकोपचाराद् ग्रहणसिद्धिः}

लोकोपचारात् जनव्यवहारात् पूर्वाचार्यशास्त्रात् अनुक्तस्य ग्रहणस्य कार्यस्य सिद्धिर् निष्पत्तिः भवति। ये लुग्लोपप्लतादयः ये इह शास्त्रे अनुक्ताः ते अन्यशास्त्राद् अनुसारणीयाः।

इति लघुवृत्तौ सन्धि प्रकणे संज्ञापादः प्रथमः॥१॥

\subsection{समानपादः}

\sutra{समानः सवर्णे दीर्घीभवति परश्च लोपम्}

समानवर्णः सवर्णे परे दीर्घवान् भवति परश् वर्णः लोपम् आपद्यते। माषाढकम् दधीहते मधूहनं होत्खषिः सम्ल्हकारः।

\sutra{अवर्ण इवर्णे ए}

अवर्णः अकाराकौ इवर्णे इकारे ईकारे च परे एकारो भवति। परलोपश्च तवेदं सेयम् उमेहते।।

\sutra{उवर्णे ओ}

अवर्णः अकाराकारौ उवर्णे उकारे ऊकारे च परे ओकारो भवति परलोपश्च। तवोदकं गङ्गोदकम् उमोहते।।

\sutra{ऋवर्णे अर्}

अवर्णः अकाराकारौ ऋवर्णे ऋकारे ॠकारे च परे अर् भवति परलोपश्च तवर्णसिद्धिः तवर्कारः।

\sutra{ऌवर्णे अल्}

अवर्णः अकाराकारौ ऌकारे ॡकारे च परे अल् भवति परलोपश्च। तवल्कारः। खट्वल्कारः।।

\sutra{एकारे ऐ ऐकारे च}

अवर्णः अकाराकारौ एकारे ऐकारे च परे ऐकारो भवति परलोपश्च। तवैडिका खट्वैतिकायनम् उमैषा।।

\sutra{ओकारे औ औकारे च}

अवर्णः अकाराकारौ ओकारे औकारे च परे औकारो भवति परलोपश्च। तवौदनः खट्वौपगवः।।

\sutra{इवर्णो यम् असवर्णे न च परलोपः}

इवर्णः इकारईकारौ असवर्णे स्वरे परे यत्वमापद्यते परश्च वर्णो लोपं नापद्यते। दध्यत्र नद्येषा।।

\sutra{वम् उवर्णः}

उवर्णः उकारऊकारौ असवर्णे स्वरे परे वत्वम् आपद्यते परश्च वर्णो लोपं नापद्यते। मध्वत्र। वध्वाख्या।

\sutra{रमृवर्णः}

ऋवर्णः ऋकारऋकारौ असवर्णे स्वरे परे रत्वम् आपद्यते परश्च वर्णो लोपं नापद्यते। धिन्नर्थः मात्रर्थः राकृतिः।।

\sutra{लमॢवर्णः}

 ऌवर्णः ऌकारॡकारौ असवर्णे स्वरे परे लत्वम् आपद्यते परश्च लोपं नापद्यते। लनुबन्धः लाकृतिः।।

\sutra{ए अय्}

एकारः अय् भवति स्वरे परे। अग्नयः बुद्धयः।।

\sutra{ऐ आय्}

ऐकारः आय् भवति स्वरे परे रायौरायः।।

\sutra{ओ अव्}

ओकारः अव् भवति स्वरे परे। धेनवः पटवः।।

\sutra{औ आव्}

औकारः आव् भवति स्वरे परे नावौनावः।।

\sutra{अयादीनां यवलोपः पदान्ते न वा लोपे तु प्रकृतिः}

अय् आय् अव् आव् एते अयादयः एषाम् अयादीनां पदान्ते स्थितानां यौ यकारवकारौ तयोर् लोपो भवति विकल्पेन। लोपे तु कृते सति प्रकृतिः यथास्वरूपं भवति पुनः सन्धिकार्यं न प्राप्नोतीत्यर्थः। कग्नासतेवा कयासते कस्मा इदं वा कस्मायिदं। इन्द इति वा इन्द विति। असा इह वा असाविह।।

\sutra{एदोत्परः पदान्ते लोपमकारः}

पदान्ते यौ एदो तौ तयोः परः अकारः लोपम् आपद्यते। अग्रेsत्र वायोsत्र।।

\sutra{न व्यञ्जने स्वराः सन्धेयाः}

व्यञ्जने परे स्वराः सन्धेयाः सन्धातव्या न भवन्ति। दण्डहस्तः दधिशीतम् मधुभाजतम् गौरीगृहम् मातृमण्डलम्।।

इति लघुवृत्तौसन्धिप्रकरणे समानपादो द्वितीयः।।२।।

\subsection{ओदन्तपादः}

\sutra{ओदन्ता अइउआनिपातः स्वरे प्रकृत्या}

ओदन्ता ओकारान्तानिपाताः अइउआश्च स्वरे परे प्रकृत्या भवन्ति सन्धिकृतां विकृतिं नापद्यन्ते इत्यर्थः। नो एतत् अहो अद्याः हंहोअधमाः अअपेहिइइहास्स्वउउत्तिष्ठआएवं नुमन्यमे आएवं नुकिलतत्।।

\sutra{द्विवचनमनौ}

अनौ औकाररहितं द्विवचनं प्रकृत्या भवति। अग्री अत्र वायू अत्र अनौ इति किम् तावतन्न।।

\sutra{बहुवचनमी}

अमी इत्येतद् बहुवचनं प्रकृत्या भवति स्वरे परे अमी आढ्याः बहुवचनमिति किम् अस्यत्र।।

\sutra{अनुपदिष्टाश्च}

दीर्घैः सूचिता वर्णसमाम्नाये साक्षाद् ये नोक्ताः प्लुताः ते स्वरे परे प्रकृत्या भवन्ति। चैत्रा इह चैत्रा आगच्छ।
एकमात्रो भवेद् ह्रस्वो द्विमात्रो दीर्घ उच्यते।
त्रिमात्रस्तु प्लुतो ज्ञेयो व्यञ्जनं त्वर्धमात्रकम्।।

इति लघुवृत्तौ सन्धिप्रकरणे ओदन्तपादस् तृतीयः॥३॥

\subsection{वर्गपादः}

\sutra{वर्गप्रथमाः पदान्ताः स्वरघोषवत्सु तृतीयान्}

वर्णाणां ये प्रथमाः कचटतपाः पदान्तास् ते स्ववर्गात् प्रत्यासन्तेस् तृतीयान् आपद्यन्ते स्वरेषु परेषु घोघवत्सु च। वागत्र अजिह षडानय तदिह त्रिष्टुबियम्। वाग्गतिः अज्गतिः षड्गच्छन्ति तद्गमनम् त्रिष्टुब्गायत्री।।

\sutra{पञ्चमे पञ्चमांस् तृतीयान् वा}

वर्गप्रथमाः पदान्ताः स्ववर्गात् पञ्चमे परे पञ्चमान् आपद्यन्ते तृतीयान् वा। वाङ्मधुरा वा वाग्मधुरा। अञ्मात्रं वा अज्मात्रम्। षड्नयनं वा षण्नयनम्। तद्नीतिः वा तन्नीतिः। त्रिष्टुम्मिनोति वा त्रिष्टुब्मिनोति।।

\sutra{वर्गप्रथमेभ्यः शकारः स्वरयवरपरच्छकारं न वा}

वर्गप्रथमेभ्य पदान्तेभ्यः परः शकारः स्वरपरः यपरः वपरः रेफपरश्च छकारम् आपद्यन्ते। नवा विकल्पेन। 
वाक्छूरः वा वाक्शूरः। अच्छ्यासः वा अच्श्यासः। षट्छ्यामा वा षट्श्यामा। तच्छेतं वा तच्शेतं। त्रिष्टुप्छ्रूयते वा त्रिष्टुप्श्रूयते।।

\sutra{तेभ्य एव हकारः पूर्वचतुर्थं न वा}

तेभ्यः वर्गप्रथमेभ्यः पदान्तेभ्यः परः हकारः पूर्वस्य वर्गस्य सम्बन्धिनं चतुर्थम् आपद्यते न वा विकल्पेन। वाग्घीना वा वाग्हीना। अज्झलानि वा अज्हलानि। षड्फलानि वा षड्हलानि। तद्धितं वा तद्हितम्। त्रिष्टुब्भुतम् वा त्रिष्टुब्हुतम्।।

\sutra{पररूपं तकारो लचटवर्गेषु}

तकारः पदान्तः लकारे चवर्गे टवर्गे च परे पररूपम् आपद्यते। तल्लुनाति। तच्चिनोति। तच्छामयति। तज्जयति। तज्झषति। तञ्ञकारः। तट्टीते। तट्xकूरः तड्डीनः। तड्ढौकते। तल्लकारः।।

\sutra{चं शे}

तकारः पदान्तः शे शकारे परे चत्वमापद्यते। तच्श्लाघ्यम्। तच्श्मशानम्।

\sutra{ङणना ह्रस्वोपधाः स्वरे द्विः}

ङणनाः पदान्ताः ह्रस्वोपधाः स्वरे परे द्विर् भवन्ति। उदङ्ङत्र। सुगण्णाह। पवन्निह।

\sutra{नोsन्तश्चछयोः शकारम् अनुस्वारपूर्वम्}

नकारः पदान्तः चछयोः परयोः अनुस्वारपूर्वं शकारम् आपद्यते। भवांश्चिनोति। भवांश्छादयति।

\sutra{टठयोः षकारम्}

नकारः पदान्तः टठयोः परयोः अनुस्वारपूर्वं षकारम् आपद्यते। भवांष्टीकते। भवांष्ठकुरः।

\sutra{तथयोः सकारम्}

नकारः पदान्तः तथयोः परयोः अनुस्वारपूर्वं सकारम् आपद्यते। भवांस्तनोति। भवांस्थुडति।

\sutra{ले लम्}

नकारः पदान्तः ले लकारे परे अनुस्वारपूर्वं लकारम् आपद्यते। भवांल्लुनाति। भवांल्लिखति।

\sutra{जझञशकारेषु ञकारम्}

नकारः पदान्तः जकारे झकारे ञकारे शकारे च परे ञकारम् आपद्यते। भवाञ्जयति। भवाञ्जकारः। भवाञ्शेते।

\sutra{शि ञ्चं वा}

नकारः पदान्तः शि शकारे परे ञ्चत्वम् आपद्यते वा विकल्पेन। महाञ्च्छूरः। महाञ्च्शूरः। महाञ्शूरः।

\sutra{डढणपरस्तु णकारम्}

नकारः पदान्तः डकारे ढकारे णकारे च परे णकारम् आपद्यते। महाण्डीनः। महाण्ढौकते। महाण्णकारः।

\sutra{मोsनुस्वारं व्यञ्जने}

मकारः पदान्तः अनुस्वारम् आपद्यते व्यञ्जने परे। त्वं हससि। द्वं रमसे।

\sutra{वर्ग्ये तद्वर्गपञ्चमं वा}

मकारः पदान्तः वर्ग्ये वर्गाक्षरे परे तस्य वर्ग्यस्य यः वर्गः तस्य पञ्चमं वा आपद्यते। 
त्वङ्करोषि। त्वंकरोषि। त्वञ्चरसि वा त्वंचरसि। त्वण्टीकसे वा त्वंटीकसे। त्वन्तरसि वा त्वंतसि। त्वम्पासि वा त्वंपासि।

इति लघुवृत्तौ सन्धिप्रकरणे वर्गपादश् चतुर्थः॥४॥

\subsection{विसर्जनीयपादः}

\sutra{विसर्जनीयश् चे छे वा शम्}

विसर्जनीयः पदान्तः चे वा छे वा परे शकारम् आपद्यते। कश्चिनोति। कश्छादयति।

\sutra{टे ढे वा षम्}

विसर्जनीयः पदान्तः टे वा ढे वा परे षत्वम् आपद्यते। कष्टीकते। कष्ठकुरः।

\sutra{ते थे वा सम्}

विसर्जनीयः पदान्तः ते वा थे वा परे सत्वम् आपद्यते। कस्तनोति। कस्थुडति।

\sutra{कखयोर् जिह्वामूलीयं न वा}

विसर्जनीयः पदान्तः कखयोः परयोः जिह्वामूलीयम् आपद्यते न वा विकल्पेन। कᳳकरोति वा कःकरोति। कᳳखनति वा कःखनति।

\sutra{पफयोर् उपध्मानीयं न वा}

विसर्जनीयः पदान्तः पफयोः परयोः उपध्मानीयम् आपद्यते न वा विकल्पेन। कᳲपाति वा कःपाति। कᳲफलति वा कःफलति।

\sutra{शे षे से वा वा पररूम्}

विसर्जनीयः पदान्तः शे वा षे वा से वा परे पररूपम् आपद्यते वा विकल्पेन।  कश्शेते वा कःशेते। कष्षष्ठः वा कःषष्ठः। कस्साधुः वा कःसाधुः। 

\sutra{उम् अकारयोर् मध्ये}

द्वयोर् अकारयोः मध्ये यः विसर्जनीयः स उत्वम् आपद्यते। कोsर्थः। सोsत्र।

\sutra{अघोषवतोश्च}

अकारघोषवतोः मध्ये यः विसर्जनीयः स उत्वम् आपद्यते। कोगतः। योधावति।

\sutra{अपरो लोप्यो sन्यस्वरे यं वा}

अकारात् परो यो विसर्जनीयः सो sन्यस्मिन्नाकारादौ स्वरे परे लोप्यो लोपनीयो भवति वा विकल्पेन यत्वम् आपद्यते। कआह वा कयाह। सइह वा सयिह। रुद्रउपरि वा रुद्रयुपरि। 

\sutra{आभोभ्याम् एवमेव स्वरे}

आकारभोशब्दाभ्यां परो विसर्जनीयः लोप्यो भवति यत्वं वापद्यते सर्वस्मिन् स्वरे परे। देवाआहुः देवायाहुः। भोइदं वा भोयिदम्। 

\sutra{घोषवति लोपम्}

आकारभोशब्दाभ्यां परो विसर्जनीयः लोपं प्राप्नोति घोषवति परे। देवायान्ति। भोयज।

\sutra{नामिपरो रम्}

नामिनः परो विसर्जनीयः रत्वं प्राप्नोति। सुपीः। सुतूः।

\sutra{घोषवत्स्वरपरः}

नामिनः परो वसर्जनीयः रत्वम् आपद्यते घोषवति स्वरे परे। अग्निर्गतः। अग्रिरत्र। पटुर्गतः। पटुरत्र।

\sutra{रप्रकृतिर् अनामिपरोपि}

रःप्रकृतिः रेफप्रकृतिः रेफस्य स्थाने विहितो यः विसर्जनीयः सः अनामिपरोपि रत्वं प्राप्नोति घोषवति स्वरे परे। अन्तर्गतः। अन्तरत्र।

\sutra{एषसपरो लोप्यो व्यञ्जने}

एषसशब्दाभ्यां परः विसर्जनीय- लोप्यो भवति व्यञ्जने परे। एषचरति। सचरति। एषददाति। सददाति।

\sutra{न विसर्जनीयलोपे पुनः सन्धिः}

विसर्जनीयस्य लोपे कृते सति पुनः सन्धिकार्यं प्राप्तं न भवति। कआह। देवाआहुः। भोइदम्।

\sutra{रो रे लोपं स्वरश्च पूर्वो दीर्घः}

रः रेफः रेफेः परे लोपं प्राप्नोति पूर्वः स्वरश्च दीर्घो भवति। अग्रीरक्षति। इन्दूरमते।

\sutra{द्विर्भावं स्वरपरश् छकारः}

स्वरात् परः छकारः द्विर् भावं द्वित्वं प्राप्नोति। वृक्षच्छाया। देवच्छत्रम्।

इति लघुवृत्तौ सन्धिप्रकरणे विसर्जनीयपादः पञ्चमः॥५॥

\subsection{निपातपादः}

\sutra{चवाहाहैवनोहंहोउताहोअहोअथोनिपातोsव्ययम्}

च वा हा अह एव नो हंहो उता हो अहो अथो एषामेकादशानां चादीनां शब्दानां एकैको निपातसंज्ञो भवति। अव्ययसंज्ञश् च। प्रयोजनं क्व? ओदन्ताअइउअनिपाताः स्वरेप्रकृत्या अव्ययाश्चेत्यादिषु।

\sutra{अन्तः स्वः प्रातः पुनर् अधोsप्यहः}

अन्तर् स्वर् प्रातर् पुनर् अधर् अहर् एते अन्तरादयः षट् अव्ययसंज्ञा भवन्ति। अपिशब्दः समुच्चयार्थः। अन्तः रम्यम्। अन्तः तिष्ठति। अन्तरागतः। अन्तःपश्यति।

\sutra{प्रपरानिनिरुद्दुस्संव्यवानुपर्यभ्यधिप्रतिस्वाङत्यप्यपोपोपसर्गाः}

प्र परा नि निर् उत् दुर् सं वि अव अनु परि अभि अधि प्रति सु आङ् अति अपि अप उप एते विंशतिरुपसर्गसंज्ञा भवन्ति अव्ययाश्च। प्रयोजनं उपसर्गेरुवइत्यादिषु।

\sutra{अन्वाद्याः कर्मप्रवचनीयाः}

अनु परि अभि अधि प्रति सु अङ् अति अपि अप उप एते एकादशकर्मप्रवचनीयसंज्ञा भवन्ति। कर्वप्रवचनीयसंज्ञया क्व प्रयोजनम्? कर्म प्रवचनीयैश् चेत्यादिषु वृक्षमन्तविद्योतते विद्त् वृक्षंपरि वृक्षधि। 

\sutra{परकार्ये विकार्यादिः}

परस्य प्रत्ययस्य कार्ये कर्तव्ये सति आद्योवर्णोविकारीभवति आसीनः।

\sutra{पूर्वकार्येन्त्य इष्यते}

पूर्वस्य प्रकृतेः कार्ये कर्तव्ये सति अन्त्यो वर्णो विकारी भवति। सः यः।

\sutra{क्वचित्पूर्वपरौसर्वौ}

पूर्वपरौ प्रकृतिप्रत्ययौ सर्वौ विकारिणौ भवतः क्वचित् न सर्वत्र। इतः इह।

\sutra{संसार्यंयोगसाधनम्}

योगस्य सूत्रवाक्यस्य साधनं पूरणं यत् उपादेयम् तद् अन्यसूत्रेभ्यः संसार्यमनुवत्तनीयम्। यथा जसीत्यस्य योगस्य पूरणं अकारोदीर्घं घोषवतीतिसूत्राद् अनुवर्तते इति परिभाष्यते।

इति लघुवृत्तौ सन्धिप्रकरणे निपातपादः षष्ठः॥६॥

सन्धिप्रकरणंसमाप्तम्॥॥


\section{नामप्रकरणम्}

\subsection{लिङ्गपादः}

\sutra{धातुविभक्तिवर्जम् अर्थवल् लिङ्गम्}

अर्थवत् सार्थकं शब्दरूपं लिङ्गसंज्ञं भवति धातुविभक्ती वर्जयित्वा। लिङ्गसंज्ञया क्व प्रयोजनम्? लिङ्गान्तनकारस्येत्यादिषु। राजा।

\sutra{तस्मात् परा विभक्तयः।}

सि औ जस्\\
अम् औ शम्\\
टा भ्याम् भिस्\\
ङे भ्याम् भ्यस्\\
ङसि भ्याम् भ्यस्\\
ङस् ओस् आम्\\
ङि ओस् सुप्।

तस्मात् लिङ्गात् पराः स्यादयः सप्तविभक्तयो भवन्ति।

सि औ जस् इति प्रथमा।\\
अम् औ शस् इति द्वितीया।\\
टा भ्याम् भिस् इति तृतीया। \\
ङे भ्याम् भ्यस् इति चतुर्थी।\\
ङसि भ्याम् भ्यस् इति पञ्चमी।\\
ङस् ओस् आम् इति षष्ठी।\\
ङि ओस् सुप् इति सप्तमी।\\
आमन्त्रणे प्रथमैवाष्टमी।

वृक्षः वृक्षौ वृक्षाः।
वृक्षम् वृक्षौ वृक्षान्।
वृक्षेण वृक्षाभ्याम् वृक्षैः।
वृक्षाय वृक्षाभ्याम् वृक्षेभ्यः।
वृक्षात् वृक्षाभ्याम् वृक्षेभ्यः।
वृक्षस्य वृक्षयोः वृक्षाणाम्।
वृक्षे वृक्षयोः वृक्षेषु।
आमन्त्रणे हे वृक्ष हे वृक्षौ हे वृक्षाः।

विभक्तिसंज्ञया क्व प्रयोजनम्? एषां विभक्तावन्तलोप इत्यादिषु।

\sutra{पञ्चादौ घुट्}

स्यादीनाम् आदौ ये पञ्चप्रत्ययाः ते घुट् संज्ञा भवन्ति। घुट् ५ – सि औ जस् अम् औ।
घुट्संज्ञया क्व प्रयोजनम्? घुटिचासम्बुद्धावित्यादिषु। 
राजा राजानौ राजानः राजानं राजानौ पञ्चेति। किं राज्ञः पश्य।

\sutra{जस्शसौ नपुंसके}

नपुंसकलिङ्गे जस् शस् इत्येतौ घुट्संज्ञौ भवतः। घुट् २ – जस् शस्।
घुट्संज्ञया क्व प्रयोजनम्? घुटिचासस्बुद्धावित्यादिषु।
सामानि रम्याणि। सामानि पठन्ति।

\sutra{आमन्त्रिते सिः सम्बुद्धिः}

आमन्त्रितम् आह्वानम् आमन्त्रितेsर्थे यः सिप्रत्ययः स सम्बुद्धिसंज्ञो भवति। 
सम्बुद्धिसंज्ञया क्व प्रयोजनम्? सम्बुद्धौ चेत्यादिषु।
हे गङ्गे। हे यमुने।

\sutra{आगम उदनुबन्धः स्वराद् अन्त्यात् परः}

उदनुबन्धः उकारानुबन्धः आगमः अन्त्यात् स्वरात् परो भवति। अनड्वान्।

\sutra{तृतीयादौ तु परादिः}

उदनुबन्धः आगमः तृतीयादौ विभक्तिगणे परे परस्य प्रत्ययस्यादिर् अवयवो भवति सर्वेषाम्।

\sutra{इदुदग्निः}

इकारान्तः शब्दः उकारान्तः शब्दश्च अग्निसंज्ञो भवति।
अग्निसंज्ञया क्व प्रयोजनम्? अग्नेरमोकार इत्यादुषु। 
अग्निं। पटुं। बुद्धिं। धेनुम्।

\sutra{ईदूत् स्त्र्याख्यौ नदी}

ईकारान्तः ऊकारान्तश्च शब्दः स्त्र्याख्यौ स्त्रीवाचकौ नदीसंज्ञौ भवतः।
नदीसंज्ञया क्व प्रयोजनम्? नद्याऐआसासामित्यादिषु। 
नद्यै। नद्याः। नद्याः। नद्यां। वध्वै। वध्वाः। वध्वाः। वध्वाम्।

\sutra{आ श्रद्धा}

आकारान्तः शब्दः स्त्रीवाचकः श्रद्धासंज्ञो भवति।
श्रद्धासंज्ञया क्व प्रयोजनम्? श्रद्धायाः सिलोपमित्यादिषु।
गङ्गा। या। सा। का।

\sutra{अन्त्यात् पूर्व उपधा}

अन्त्याद् वर्णात् पूर्वो वर्ण उपधासंज्ञो भवति। 
उपधासंज्ञया क्व प्रयोजनम्? नान्तस्य चोपधाया इत्यादिषु।
पञ्चानाम्। सप्तानाम्। दशानाम्।

\sutra{व्यञ्जनान्नोनुषङ्गः}

अन्त्याद् व्यञ्जनात् पूर्वो नकारः अनुषङ्गसंज्ञो भवति।
अनुषङ्गसंज्ञया क्व प्रयोजनम्? अनुषङ्गश्चाकुञ्चेरित्यादिषु।
विदुषः। महतः।

\sutra{धुड् व्यञ्जनम् अनन्तस्थानुनासिकम्}

व्यञ्जनम् धुट्संज्ञं भवति अन्तःस्थानुनासिकान्वर्जयित्वा।
धुट् २४ – कखगघ चछजझ टठडढ तथदध पफबभ शषसह।
धुट्संज्ञया क्व प्रयोजनम्? अटान्तृतीयइत्यादिषु।
अग्निमध्यां विरुद्धिः।

\sutra{अकारो दीर्घं घोषवति}

अकारः लिङ्गान्तः दीर्घं प्राप्नोति घोषवदादौ विभक्तिगणे परे।
वृक्षाय। देवाभ्याम्।

\sutra{जसि}

अकारः लिङ्गान्तः दीर्घं प्राप्नोति जसि परे।
वृक्षाः। देवाः।

\sutra{शसि सस्य च नः}

अकारो लिङ्गान्तः दीर्घं प्राप्नोति शसि परे शसश्च सकारस्य नत्वं भवति।
वृक्षान्। देवान्।

\sutra{अकारे लोपम्}

अकारो लिङ्गान्तः लोपं प्राप्नोति विभक्तिसम्बन्धिन्यकारे परे।
वृक्षम्। देवम्।

\sutra{भिस् ऐस् वा}

अकाराल् लिङ्गान्तात् परो भिस्प्रत्ययः ऐसादेशं प्राप्नोति।
वृक्षैः देवैः।

\sutra{धुटि बहुत्वे त्वे}

अकारो लिङ्गान्तः एत्वं प्राप्नोति बहुत्वे उत्पन्ने धुटि परे।
वृक्षेभ्यः। देवेषु।

\sutra{ओसि च}

अकारो लिङ्गान्तः एत्वं प्राप्नोति ओसि परे।
वृक्षयोः। देवयोः।

\sutra{ङसिर् आत्}

अकारान्ताल् लिङ्गात् परो ङसिप्रत्ययः आदादेशं प्राप्नोति।
वृक्षात्। देवात्।

\sutra{ङस् स्य}

अकारान्ताल् लिङ्गात् परः ङस्प्रत्ययः स्य संपद्यते।
वृक्षस्य। देवस्य।

\sutra{इन टा}

अकारान्ताल् लिङ्गात् परः टाप्रत्ययः इनादेशं प्राप्नोति।
वृक्षेण। देवेन।

\sutra{ङेर् यः}

अकारान्ताल् लिङ्गात् परः ङेप्रत्ययः यादेशं प्राप्नोति।
वृक्षाय। देवाय।

\sutra{स्मै सर्वनाम्नः}

अकारान्तात् सर्वनाम्नः परः ङेप्रत्ययः स्मैआदेशं प्राप्नोति।
सर्वस्मै। विश्वस्मै। यस्मै। तस्मै। कस्मै। अमुष्मै। अस्मै।

\sutra{ङसिः स्मात्}

अकारान्तात् सर्वनाम्नः परः ङसिप्रत्ययः स्मात् संपद्यते।
सर्वस्मात्। यस्मात्। तस्मात्। कस्मात्। अमुष्मात्। अस्मात्।

\sutra{ङि स्मिन्}

अकारान्तात् सर्वनाम्नः परः ङिप्रत्ययः स्मिन् संपद्यते।
सर्वस्मिन् विश्वस्मिन्। यस्मिन्। तस्मिन्। कस्मिन्। अमुष्मिन्। अस्मिन्।

\sutra{विभाष्येते पूर्वादेः}

पूर्वादेर् गणात् परयोः ङसिङ्योः स्मात् स्मिन् आदेशौ \textbf{विभाष्येते} विकल्पेन कथ्येते।
पूर्वस्मात् वा पूर्वात्। पूर्वस्मिन् वा पूर्वे। परस्मात् वा परात्। परस्मिन् वा परे।

\sutra{सुर् आमि सर्वतः}

सर्वतः सर्वस्मात् अकारान्तात् आकारान्ताच्च सर्वनाम्नः \textbf{आमि}परे \textbf{सुरा}गमो भवति।
सर्वेषां येषां तेषां केषां अमीषां एषाम् सर्वासां यासां तासां कासां अमूषां आसाम्।

\sutra{जस् सर्वईम्}

अकारान्तात् सर्वनाम्नः परो जस् प्रत्ययः ईत्वं प्राप्नोति।
सर्वे ये ते के अमी इमे।

\sutra{अल्पादेर् वा}

अल्पादेर् गणात् परः जस् प्रत्ययः ईत्वं वा प्राप्नोति। 
अल्पे वा अल्पाः प्रथमे वा प्रथमाः चरमे वा चरमाः कतमे वा कतमाः।

\sutra{पूर्वादेश्च}

पूर्वादेर् गणात् परः जस्प्रत्ययः ईत्वं वा प्राप्नोति।
पूर्वे वा पूर्वाः। परे वा पराः।

\sutra{द्वन्द्वस्थाच्च}

द्वन्द्वसमासस्थात् सर्वनाम्नः परो जस्प्रत्ययः ईत्वं वा प्राप्नोति।
कतरे च कतमे च ते कतरकतमे वा कतरकतमाः। दत्तकतरकतमे वा दत्तकतरकतमाः। 

\sutra{नान्यत् सर्वनामिकम्}

द्वन्द्वसमासस्थात् सर्वनाम्नः परस्य प्रत्ययस्यान्यत्सार्वनामिकं कार्यं न भवति। 
पूर्वे चापरे च पूर्वापरं। तस्मै पूर्वापराय पूर्वापरात् पूर्वापराणां पूर्वापरे।

\sutra{तृतीया समासे}

तृतीयासमासे सति सर्वनाम्नः परस्य प्रत्ययस्यान्यत्सार्वनामिकं कार्यं न भवति।
मासेन पूर्वः मासपूर्वः। तस्मै समासपूर्वाय मासपूर्वात् सामपूर्वाणां मासपूर्वे।

\sutra{बहुव्रीहौ च}

बहुव्रीहिसमासे सति सर्वनाम्नः परस्य प्रत्ययस्य सार्वनामिकं कार्यं न भवति।
प्रियाः सर्वे यस्य सः प्रियसर्वः। तस्मै प्रियसर्वाय प्रियसर्वात् प्रियसर्वाणां प्रियसर्वे।

\sutra{दिशां वा}

दिशां दिग्वाचिनां शब्दानां यो बहुव्रीहिस् तत्रस्थितात् सर्वनाम्नः परस्य प्रत्ययस्य सार्वनामिकं कार्यं वा भवति। दक्षिणस्याश्च पूर्वस्याश्च दिशोर्यदन्तरालं विदिक् सा दक्षिणपूर्वा। तस्यै दक्षिणपूर्वस्यै वा दक्षणपूर्वायै। दक्षिणापूर्वस्याः वा दक्षिपूर्वायाः। दक्षिणपूर्वस्याः वा दक्षिणपूर्वायाः। दक्षिणपूर्वस्यां वा दक्षिणपूर्वायाम्।

\sutra{श्रद्धायाः सिः लोपम्}

श्रद्धासंज्ञकात् परः सिः प्रत्ययः लोपं प्राप्नोति।
गङ्गा या सा का।

\sutra{टौसोर् ए}

श्रद्धाया अन्त्यस्य एत्वं भवति टाओस् इत्येतयोः परयोः।
गङ्गया गङ्गयोः उमया उमयोः।

\sutra{सम्बुद्धौ च}

श्रद्धाया अन्त्यस्य एत्वं भवति सम्बुद्धौ परस्याम्।
हे गङ्गे हे युमने।

\sutra{ह्रस्वोऽम्बार्थानाम्}

अम्बार्था मातृवाचका ये श्रद्धासंज्ञकाः शब्दास्तेषाम् अन्त्यस्य ह्रस्वो भवति सम्बुद्धौ परस्याम्।
हे अम्ब हे अक्क हे अत्त हे अल्ल।

\sutra{औरीम्}

श्रद्धासंज्ञकात् परः औप्रत्ययः ईत्वं प्राप्नोति।
शाले माले।

\sutra{ङवन्ति यै यास् यास् याम्}

श्रद्धासंज्ञकाल्लिङ्गात्पराणि ङवन्तिवचनानि ङे ङसि ङस् ङि इत्येतानि यथासंख्यं यै यास् यास् याम् इत्येतानादेशान् आपद्यन्ति।\\
गङ्गायै गङ्गायाः गङ्गायाः गङ्गायाम्।\\
उमायै उमायाः उमायाः उमायाम्।

\sutra{सर्वनाम्नस् तु ससवो ह्रस्वपूर्वाः}

श्रद्धासंज्ञकात् सर्वनाम्नः परेषां ङवतां वचनानां यथासङ्ख्यम् यै यास् यास् याम् इत्येते आदेशाः ससवः स्वागमयुक्ताः ह्नस्वपूर्वाश्च भवन्ति।
सर्वस्यै सर्वस्याः सर्वस्याः सर्वस्यां यस्यै यस्याः यस्याः यस्यां तस्यै तस्याः तस्याः तस्यां कस्मै कस्याः कस्याः कस्याम् अमुष्यै अमुष्याः अमुष्याः अमुष्याम् अस्यै अस्याः अस्याः अस्याम्।

\sutra{द्वितीयातृतीयाभ्यां वा}

द्वितीयातृतीयाशब्दाभ्यां परेषां ङवतां वचनानाम् यथासङ्ख्यं यै यास् यास् याम् इत्येते आदेशाः ससवः स्वागमयुक्ता ह्रस्वपूर्वाश्च भवन्ति।
द्वितीयस्यै वा द्वितीयायै। द्वितीयस्याः वा द्वितीयायाः। द्वितीयस्याः वा द्वितीयायाः। द्वितीयस्यां वा द्वितीयायाम्। एवं तृतदीयस्यै वा तृतीयायै। तृतीयस्याः वा तृतीयायाः। तृतीयस्याः वा तृतीयायाः। तृतीयस्यां वा तृतीयायाम्।

\sutra{नद्या ऐ आस् आस् आम्}

नदीसंज्ञकाल् लिङ्गात् परेषां ङवतां वचनानां यथासङ्ख्यं ऐ आस् आस् आम् इत्येते आदेशा भवन्ति।
नद्यै नद्याः नद्याः नद्यां वध्वै वध्वाः वध्वाः वध्वाम्।

\sutra{सम्बुद्धौ ह्रस्वः}

नदीसंज्ञकाल् लिङ्गात् अन्त्यस्य ह्रस्वो भवति सम्बुद्धौ पर्याम्। 
हे नदि हे वधु।

\sutra{अम्शसोर् आदिर् लोपम्}

नदीसंज्ञकाल् लिङ्गात् परयोः अम् शस् इत्येतयोः प्रत्ययोः आदिः लोपं प्राप्नोति।
नदीं नदीः वधूं वधूः।

\sutra{ईकारान्तात् सिः}

ईकारान्तान् नदीसंज्ञकाल् लिङ्गात् परः सिप्रत्ययो लोपं प्राप्नोति।
नदी मही गौरी।

\sutra{व्यञ्जनाच्च}

व्यञ्जनान्ताल् लिङ्गात् परः सिप्रत्ययः लोपं प्राप्नोति।
वाक् विद्युत् राजा।

\sutra{अग्नेरमोऽकारः}

अग्निसंज्ञकाल् लिङ्गात् परस्य अमो sकारः लोपं प्राप्नोति। 
अग्निं पटुं बुद्धिं धेनुम्।

\sutra{औकारः पूर्वम्}

अग्निसंज्ञकाल् लिङ्गात् परः औकारः पूर्वं वर्णं प्राप्नोति।
अग्नी पटू बुद्धी धेनू।

\sutra{शसोऽकारः सचनो ऽस्त्रियाम्}

अग्निसंज्ञकाल् लिङ्गात् परस्य शसः अकारः पूर्ववर्णं प्राप्नोति। शसश् च सकारस्य च नत्वं भवति। अस्त्रियां स्त्रीविषयं वर्जयित्वा।
अग्नीन् पटून्।
अस्त्रियां किं? बुद्धीः धेनूः।

\sutra{टा ना}

अग्निसंज्ञकाल् लिङ्गात् परः टाप्रत्ययः ना आदेशं प्राप्नोति अस्त्रियां स्त्रीविषयं वर्जयित्वा।
अग्निना पटुना।
अस्त्रियामिति किम्? बुध्या धेन्वा।

\sutra{अदो ऽमुश्च}

अदश्शब्दात् परः टाप्रत्ययः ना आदेशं प्राप्नोति। 
अदसश्च अमुआदेशः भवति अस्त्रियां स्त्रीविषयं वर्जयित्वा। 
अमुना।
अस्त्रियामिति किं? अमुया।

\sutra{इरेदुरोज्जसि}

अग्निसंज्ञकाल् लिङ्गात् इः एद्भवति उः ओद्भवति जसि परे। 
अग्नयः पटवः बुद्धयः धेनवः।

\sutra{सम्बुद्धौ च}

अग्निसंज्ञकात् लिङ्गात् इरेद्भवति उरोद्भवति सम्बुद्धौ परस्याम्।
हे अग्ने हे पटो हे बुद्धे हे धेनो।

\sutra{ङे च}

अग्निसंज्ञकाल् लिङ्गात् इः एद् भवति उः ओद् भवति ङे प्रत्यये परे।
अग्नये पटवे बुद्धये धेनवे।

\sutra{ङसिङसोर् लोपश्च}

अग्निसंज्ञकाल् लिङ्गात् इः एद् भवति उः ओद् भवति ङसिङसोः परयोः तयोश्च अकारो लोपो भवति। 
अग्नेः अग्नेः पटोः पटोः बुद्धेः बुद्धेः धेनोः धेनोः।

\sutra{गोश्च}

गोशब्दात् परयोः ङसिङसोः अकारो लोपो भवति।
गोः गोः। द्योशब्दात् अपीष्यते। द्योः द्योः।

\sutra{ङिरौ सपूर्वः}

अग्निसंज्ञकाल् लिङ्गात् परो ङिप्रत्ययः पूर्वेण वर्णेन सह औकारं प्राप्नोति। 
अग्नौ पटौ बुद्धौ धेनौ।

\sutra{सखिपत्योर् ङिः}

सखिपतिशब्दाभ्यां परो ङिप्रत्ययः औत्वं प्राप्नोति।
सख्यौ पत्यौ।

\sutra{ङसिङसोरुमः}

सखिपतिशब्दाभ्यां परयोः ङसिङसोः अकार उत्वं प्राप्नोति।
सख्युः सख्युः पत्युः पत्युः।

\sutra{ऋदन्तात् सपूर्वः}

ऋदन्ताल् लिङ्गात् परयोर् ङसिङसोर् अकारः पूर्वेण वर्णेन सह उत्वं भवति। 
पितुः पितुः मातुः मातुः भ्रातुः भ्रातुः कर्तुः कर्तुः स्वसुः स्वसुः नुः नुः।

\sutra{आ सौ सिलोपश्च}

ऋदन्तस्य लिङ्गस्यान्त्यस्य आत्वं भवति सौ सिप्रत्यये परे सेश्चलोपो भवति।
पिता माता भ्राता कर्ता स्वसा ना।

\sutra{अग्निवच् छसि}

ऋदन्तं लिङ्गम् अग्निवद् भवति शसि परे। यथाग्निसंज्ञकात् परस्य शसोऽकारःसशनोऽस्त्रियामिति शसो ऽकारस्य पूर्ववर्णादेशः स्त्रीविषयं वर्जयित्वा सकारस्य च नत्वं भवति तथा ऋदन्तात् अपि भवतीत्यर्थः। 
पितॄन् भ्रातॄन् कर्तॄन् नॄन्।
स्त्रियां सस्य नत्वं नास्ति यथा मातॄः स्वसॄः।

\sutra{अर्ङौ}

ऋदन्तस्य लिङ्गस्य अन्त्यस्य अर्भवति ङौ ङिप्रत्यये परे।
पितरि मातरि भ्रातरि कर्तरि स्वसरि नरि।

\sutra{घुटि च}

ऋदन्त्यस्य लिङ्गस्यार्भवति घुटि परे।
पितरौ पितरः पितरं पितरौ मातरौ मातरः मातरं मातरौ नरौ नरः नरं नरौ।

\sutra{धातोस्तृशब्दस्यार्}

धातोः परः यस्तृशब्दस् तस्य आर्भवति चुटि परे।
कर्तारौ कर्तारः कर्तारं कर्तरौ वक्तारौ वक्तारः वक्तारं वक्तारौ।

\sutra{स्वस्रादीनां च}

स्वस्रादीनां शब्दानां यः ऋकारस् तस्य आर् भवति घुटि परे।
स्वसारौ स्वसारः स्वसारं स्वसारौ नप्तारौ नप्तारः नप्तारं नप्तारौ।
स्वसा नप्ता च नेष्टा च त्वष्टा क्षत्ता तथैव च।
होता पोता प्रशस्ता चेत्यष्टौ स्वस्रादयः स्मृताः॥

\sutra{आ च न सम्बुद्धौ}

ऋदन्तस्य लिङ्गस्य आर् न भवति। आसौसिलोपश्चेत्यनेन आत्वं न भवति। पारिशेष्याद् घटि चेत्य् अर् एव भवति सम्बुद्धौ परस्यां।
हे पितः हे मातः हे भ्रातः हे कर्तः हे स्वसः हे नः।

\sutra{ह्रस्वनदीश्रद्धाभ्यः सिर्लोपम्}

ह्रस्वान्तात् नदीसंज्ञकात् श्रद्धासंज्ञकाच्च लिङ्गात् परः सम्बुद्धिसिकारः लोपं प्राप्नोति।
हे वृक्ष हे नदि हे वधु हे गङ्गे हे यमुने।

\sutra{आमि च नुः}

ह्रस्वनदीश्रद्धाभ्यः आमि परे नुरागमो भवति।
वृक्षाणां अग्नीनां वायूनां पितॄणां नदीनां वधूनां गङ्गानाम् उमानाम्।

\sutra{त्रेस्त्रयश्च}

त्रिशब्दस्यामिपरे नुरागमो भवति। त्रेश्च त्रयादेशः। 
त्रयाणाम् परमत्रयाणाम्।

\sutra{चतुरः}

चत्वार्शब्दस्य आमि परे नुरागमो भवति।
चतुर्णाम्।

\sutra{सङ्ख्यायाः ष्णान्तायाः}

षकारनकारान्तायाः सङ्ख्याः आमि परे नुरागमो भवति।
षण्णां पञ्चानां सप्तानां दशानाम्।

\sutra{कतेश्च जस्शसोर् लुक्}

षकारनकारान्तायाः सङ्ख्यायाः कतिशब्दाश् च परयोः जस्शसोः लुग्भवति। 
षट् तिष्ठन्ति। षट् पश्य। पञ्च पञ्च सप्त सप्त नव नव। दश तिष्ठन्ति दश पश्य। कति तिष्ठन्ति कति पश्य।

\sutra{नियो ङिराम्}

णीञ् प्रापणे इत्यस्मात् क्विवन्तात् परो ङिप्रत्ययः आम् आदेशं प्राप्नोति।
नियां सेनान्यां ग्रामण्याम्।

\textbf{इति लघुवृत्तौ नाम प्रकरणे लिङ्गपादः प्रथमः॥१॥
}

\subsection{सखिपादः}

\sutra{न सखिष्टादावग्निः}

सखिशब्दः अग्निसंज्ञो न भवति। टादौ प्रत्यये परे सख्या सख्ये। टादाविति किम्? सखीन्। 

\subsection{युष्मत्पादः}

\sutra{युष्मदस्मदोः पदं पदात् षष्ठीचतुर्थीद्वितीयासु वस्नसौ}

युष्मदस्मदोः यत्पदं पदात्परं षष्ठीचतुर्थीद्वितीयासु बहुत्वे स्थितं तस्य यथासङ्ख्यं वस् नस् इत्येतावादेशौ भवतः शिवोवः स्वासीशिवोयुष्माकम् शिवोनः स्वामीशिवोस्माकम् शिवोवोदास्यतिशिवोयुष्मभ्यम् शिवोनोदास्यतिशिवोस्मभ्यम् शिवोवोरक्षतिशिवोयुष्मान् शिवोनोरक्षति शिवोस्मान्।

\subsection{कारकपादः}

\sutra{अव्ययीभावादकारान्ताद्विभक्तीनाममपञ्चम्याः}

\textbf{अकारान्तादव्ययीभावा}त्परासां \textbf{विभक्तीनाम्} \textbf{अमा}देशो भवति \textbf{अपञ्चम्याः} पञ्चमीविभक्तिं वर्जयित्वा। उपकुम्भं तिष्ठति। उपकुम्भं पश्य। उपकुम्भं देहि। उपकुम्भं स्वम्। अकारान्तादिति किम्? समीपं वध्वाः उपवधु। स्वरोह्रस्व इति ह्रस्वः। अन्यस्माल्लुगिति सिलोपः। अपञ्चम्या इति किम्? उपकुम्भादागतः।

\sutra{वा तृतीयासप्तम्योः}

अकारान्तादव्ययीभावात् परयोः तृतीयासप्तम्योः अमादेशो वा भवति। उपकुम्भं कृतम् उपकुम्भेन कृतम्। उपकुम्भं स्थापय उपकुम्भे स्थापय। अकारान्तादिति किम्? उ पवधु कृतम् उपवधु स्थापय।

\sutra{अन्यस्माल्लुक्}

\textbf{अन्यस्मात्} अनकारान्तात् अव्ययीभावात् परासां विभक्तीनां लुक् भवति। उवधु तिष्ठति। उपवधु पश्य। उपवधु कृतम्। उपवधु देहि। उपवधु आगतः। उपवधु स्वम्। उपवधु स्थापय। समृद्धिः कर्तृणां सुकर्तृ। पश्चान्नद्याः अनुनदि। विधिमनतिक्रम्य यथाविधि।

\sutra{अव्ययाच्च}

अव्ययात्परासां विभक्तीनां लुग्भवति । च वा हा अह एव ग्रामश्च ते स्वम्। नगरं च मे स्वम्।

\sutra{रूढानां बहुत्वे स्त्रियामपत्यप्रत्ययस्य}

ये शब्दाः देशस्य राज्ञां च वाचकाः ते रूढाः रूढानां सम्बन्धिनः अपत्यप्रत्ययस्य बहुत्वे उत्पन्नस्य लुग्भवति अस्त्रियां स्त्रीवषयं वर्जयित्वा। अङ्गस्यापत्यानि बहूनि अङ्गाः। एवं बङ्गाः कलिङ्गाः सुह्माः रघोरपत्यानि बहूनि रघवः। वाणपत्ये इति अण् तस्य अनेन लुक्। रूढानामिति किम्? पूरोरपत्यानि बहूनि पौरवाः अण् णकारो वृद्धिरादौसणे इत्यर्थः। आदिवृद्धिश्च उवर्णस्त्वोत्वमापाद्य इति उवर्णस्य ओकारः। बहुत्वे इति किम्? अङ्गस्य अपत्यं आङ्गः। अस्त्रियामिति किम्? अङ्गस्यापत्यानि बह्व्यः स्त्रियः आङ्ग्यः। अणेयेकादिति ईप्रत्ययः ईकारे स्त्रीकृतेल्लोप्य इति अकारलोपः जस् \textbf{अपत्यप्रत्ययस्ये}ति किम्? पञ्चालानामिमे शिष्याः पाञ्चालाः तस्येदमित्यण्।

\sutra{गर्गयस्कविदादीनां च}

गर्गादीनां यस्कादीनां विदानीनां च सम्बन्धिनः अपत्यप्रत्ययस्य बहुत्वे उत्पन्नस्य लुग्भवति अस्त्रियां स्त्रीविषयं वर्जयित्वा गर्गस्यापत्यानि बहूनि गर्गाः एवं वत्साः वजाः सङ्कृतयः व्याघ्रपदो अपत्यानि बहूनि व्याघ्रपदः। ण्यगर्गादेरितिण्यतस्यानेन लुक्। यस्कस्यापत्यानि बहूनि यस्काः लह्याः द्रूह्याः। विदस्यापत्यानि बहूनि विदाः उर्वाः भरद्वाजः वाणपत्ये इति अण्। तस्यानेन लुक्। \textbf{अपत्यप्रत्ययस्ये}ति किम्? गर्गस्य इमे शिष्याः गार्गाः यास्काः वैदाः। \textbf{बहुत्वे} इति किम्? गर्गस्यापत्यं गार्ग्यः यास्कः वैदः। अस्त्रियामिति किम्? गर्गस्यापत्यानि बह्व्यः स्त्रियाः गार्ग्यः यास्क्यः वैद्यः गर्गशब्दात् ण्यगर्गादेरितिण्यः ण्यदिति ईप्रत्ययः यस्यापत्यस्येतियलोपः शिष्टेभ्यः अण्। अणेयेकादिति ईप्रत्ययः गर्गवत्सवजसंकृतिव्याघ्रपात् अजविदभृत् इत्यादयोगर्गादयः यस्कलह्यद्रूह्यअयस्थूणबलन्धर इत्यादयो यस्कादयः विदे उर्व भरद्वाजकुर्वण उपमन्यु इत्यादयो विदादयः।

\sutra{भृग्वत्र्यङ्गिरस्कुत्सवसिष्ठगोतमेभ्यश्च}

भृगु अत्रि अङ्गिरस् कुत्स वसिष्ट गोतम एभ्यः परस्य अपत्यप्रत्ययस्य बहुत्वे उत्पन्नस्य लुक् भवति अस्त्रियां स्त्रीविषयं वर्जयित्वा। भृगोरपत्यानि बहूनि भृगवः अत्रयः अङ्गिरसः कुत्साः वसिष्ठाः गोतमाः अस्यानेन लुक्। अत्रिशब्दात् स्त्र्यात्र्यादेरेयण् अस्य अनेन लुक्। शिष्टेभ्यो वाणपत्येण्। अपत्यप्रत्ययस्येति किम्? भृगोरिमे शिष्याः भार्गवाः। ततस्येदमित्यण्। बहुत्वे इति किम्? भृगोरपत्यं भार्गवः एवम् आत्रेयः आङ्गिरसः कौत्सः वासिष्ठः गौतमः। अस्त्रियामिति किम्? भृगोरपत्यानि बह्व्यः स्त्रियः भार्गव्यः आत्रेय्यः आङ्गिरस्यः कौत्स्यः वासिष्ठ्यः गौतम्यः अणेरेकादिति ईप्रत्ययः ईकारेस्त्रीकृतेल्लोप्यः जस्।

\sutra{यतोऽपैति भयमादत्ते वा तदपादानम्}

यतः अपैति यस्माद्विश्लेषं प्राप्नोति यतो भयं यतश्चादत्ते गृह्णाति तत्कारकं अपादानसंज्ञं भवति। अपादाने पञ्चमी। ग्रामादागच्छति। अश्वात्पततिः। चौरेभ्यो बिभेति। अधर्मान्निवर्तते। कोष्ठागाराद्धान्यं गृह्णाति। उपाध्यायाद्विद्यामादत्ते। अपादानसंज्ञया क्व प्रयोजनम्? शेषाः कर्मकारणसंप्रदानापादानस्वाम्याद्यधिरकणेषु इत्यत्र वक्ष्यमाणाभिरत्रैव प्रयोजनम्।

\sutra{ईप्सितं च रक्षार्थानाम्}

रक्षार्थानां रक्षाप्रयोजनानाम् धातूनां प्रयोगे सति यदीप्सितमभिप्रीतं तत्कारकमपादानसंज्ञं भवति। च शब्दादनीप्सितमपि। यवेभ्यो गा रक्षति। तिलेभ्यः काकान्वारयति। अनीप्सितं यथा अहिभ्यः पुत्रान्रक्षति कूपादन्धं वारयति।

\sutra{यस्मै दित्सा रोचते धारयते वा तत्संप्रदानम्}

यस्मै दित्सा दातुमिच्छा यस्मै रोचते यस्मै च ऋणं धारयते तत्कारकं संप्रदानसंज्ञं भवति। संप्रदाने चतुर्थी। विप्राय गां ददाति। वैश्रवणाय बलिमुपहरति। देवदत्ताय रोचते मोदकः। यज्ञदत्ताय शतं धारयते। दित्सा इति किम्? राज्ञो दाण्डं ददाति।

\sutra{य आधारस्तदधिकरणम्}

य आधारः तत्कारकं अधिकरणसंज्ञं भवति। अधिकरणे सप्तमी। कटे आस्ते। पर्यङ्के शेति। जले मत्स्याः। खे शकुनयः दध्नि घृतम्। तिलेषु तैलम्।

\sutra{येन क्रियते तत्करणम्}

येन क्रियते साध्यते तत्कारकं करणसंज्ञं भवति। करणे तृतीया। दात्रेण लुनाति। परशुनाच्छिनत्ति।

\sutra{यत्किर्यते तत्कर्म}

यत्क्रियते क्रियया विशिष्टं संपाद्यते तत्कारकं कर्मसंज्ञं भवति। कर्मणि द्वितीया। कटं करोति। ओदनं पचति। आदित्यं पश्यति।

\sutra{यः करोति स कर्ता}

यः करोति क्रियां निर्वर्तयति स कर्तृसंज्ञो भवति। कर्तरि तृतीया। रुद्रेण त्रिपुरं दग्धम्। विष्णुना बद्धो बलिः। कर्तृसंज्ञया क्व प्रयोजनम्? कर्तरि चेत्यादिषु।

\sutra{कारयति यः स हेतुश्च}

\textbf{य} कर्तारं \textbf{कारयति} क्रियायां प्रेरयति स हेतुसंज्ञो भवति। चशब्दात्कर्तृसंज्ञश्च। कर्तरि प्रथमा। कारयति पाठयति हेतुसंज्ञया क्व प्रयोजनम् धातोश्च हेतावित्यादिषु।

\sutra{तेषां परमुभयप्राप्तौ}

तेषामपादानादीनां कारकाणां मध्यात् उभयोर्द्वयोः कारकयोः एकत्र प्राप्तौ सत्यां यत्परं तदेवकारकं भवति। वृक्षं त्यजति खगः। त्यजति दण्डं दण्डी। अत्र परत्वात्कर्मकर्तृकारके भवतः।

\sutra{प्रथमाविभक्तिर्लिङ्गार्थवचने}

लिङ्गार्थस्य वचने अभिधाने सति लिङ्गात्प्रथमाविभक्तिर्भवति। वृक्षः कुमारी। कुण्डम्। द्रोणः। आढकम्।

\sutra{आमन्त्रणे च।}

आमन्त्रणादिके लिङ्गार्थवचने सति प्रथमाविभक्तिर्भवति। हे शिव हे केशव।

\sutra{शेषाः कर्मकरणसंप्रदानापादानस्वाम्याद्यधिकरणेषु}

शेषा द्वितीयाद्याविभक्तयः कर्मणि करणे संप्रदाने अपादाने स्वाम्यादियोगे अधिकरणे च भवन्ति। कर्मणि द्वितीया। कटं करोति। करणे तृतीया। परशुनाच्छिनत्ति। संप्रदाने चतुर्थी। विप्राय गां ददाति। अपादाने पञ्चमी। ग्रमादागच्छति। अश्वात्पतितः। स्वाम्यादियोगे षष्ठी। देवदत्तस्य स्वम्। पितुः पुत्रः। राज्ञो भार्या। अधिकरणे सप्तमी। कटे आस्ते। पर्यङ्के शेते।

\sutra{पर्यपाङ्योगे पञ्चमी}

परि अप आङ् इत्येतैर्योगे सति लिङ्गात् पञ्चमी भवति। परिस्रुघ्र्नाद्वृष्टो देवः। अपकान्यकुब्जाद्वृष्टो देवः। आकुरुक्षेत्राद्वृष्टो देवः।

\sutra{दिगितरर्तेऽन्यैश्च}

दिक्शब्देरितर ऋते अन्य इत्येतैश्च योगे लिङ्गात्पञ्चमी भवति। पूर्वो ग्रामात्। दक्षिणो नगरात्। इतरो देवदत्तात्पण्डितः। ऋते विष्णोः कुतोमोक्षः। अन्यो देवदत्तात्।

\sutra{द्वितीयैनेन}

एनेन एनप्रत्ययान्तेन शब्देन योगेन सति लिङ्गाद्द्वितीया भवति। उत्तरेण ग्रामं वसति। दक्षिणेन नगरं वसति। पुरस्तादादयश्चेति एनप्रत्ययः।

\sutra{कर्मप्रवचनीयैश्च}

अन्वाद्याः एकादश कर्मप्रवचनीयाः ये प्रागुक्ताः तैः योगे सति लिङ्गाद्द्वितीया भवति। वृक्षमनुविद्योतते विद्युत्। वृक्षं परि। वृक्षमभि। मां प्रति साधुः।

\sutra{गत्यर्थकर्मणि द्वितीयाचतुर्थौ चेष्टायामनध्वनि}

गत्यर्थानां धातूनां यत्कर्म तस्मिन्कर्मणि अनध्वनि अध्वशब्दवर्जिते द्वितीयाचतुर्थ्यौ भवतः चेष्टायां गम्यमानायां सत्यां चेष्टाकायपरिस्पन्दः ग्रामं गच्छति ग्रामाय गच्छति स्रुघ्नं व्रजति। स्रुघ्नाय व्रजति। तीर्थं याति तीर्थाय याति। गत्यर्थग्रहणं किम्? ओदनं पचति कर्मणीति किम्? अश्वेन याति। चेष्टायामितिकिम्? मनसा वाराणसीं गच्छति। अनध्वनीति किम्? अध्वानं गच्छति।

\sutra{मन्यकर्मणि चानादरे प्राणिनि}

मनज्ञाने इत्यस्य यत्कर्म तस्मिन्कर्मणि अप्राणिनि प्राणिवर्जिते द्वितीयाचतुर्थ्यौ भवतः अनादरे गम्यमाने अनादरः तिरस्कारः। तृणमपि त्वा न मन्ये। तृणायापि त्वा न मन्ये। बुसमपि त्वा न मन्ये। बुसायापि त्वा न मन्ये। मन्यग्रहणे किम्? तृणं त्वां चिन्तयामि। अनादरे इति किम्? सुवर्णं त्वां मन्ये। अप्राणिनि इति किम्? न त्वा काकं मन्ये। न त्वा श्वानं मन्ये।

\sutra{नमस्स्वस्तिस्वाहास्वधालंवषड्योगे चतुर्थी}

नमस् स्वस्ति स्वाहा स्वधा अलं वषट् इत्येतैर्योगे सति लिङ्गात् चतुर्थी भवति। नमो रुद्राय। स्वस्ति प्रजाभ्यः। स्वाहा अग्नये। पितृभ्यः स्वधा। अलं मल्लो मल्लाय। वषडिन्द्राय।

\sutra{तुमर्थाच्च भाववाचिनः}

तुमासमानार्थो भाववाची घञादिप्रत्ययः भाववाचिनश्चेति अनेन क्रियायां क्रियार्थायां विहितः तदन्तात् शब्दात् चतुर्थीभवति। पाकाय व्रजति। पक्तये व्रजति। पचनाय व्रजति। पक्तुं व्रजति इत्यर्थः। तुमर्थादितिकिम्? पाकः पक्तिः पचनम्। भाववाचिन इति किम्? कटं कारको व्रजति कर्तुं व्रजति इत्यर्थः।

\sutra{तृतीया सहयोगे}

सहार्थेन शब्देन योगे सति लिङ्गात्तृतीया भवति। पुत्रेण सह आगतः। पुत्रेण सह गोमान्। पुत्रेण सत्रा। पुत्रेण सजूः।

\sutra{हेत्वर्थे}

फलसाधनयोग्यः पदार्थो हेतुः हेतावर्थे लिङ्गात्तृतीया भवति। विद्यया यशः। कन्यया शोकः। धनेन कुलम्।

\sutra{कुत्सिते ऽङ्गे}

कुत्सिते अङ्गिनि यदङ्गं तस्मिन्नङ्गे तृतीया भवति। अक्ष्णा काणः। नासिकया चिपिटः। पादेन खञ्जः। शिरसा खलतिः। पाणिना कुणिः।

\sutra{विशेषणे}

विशेष्यते उपलक्ष्यते येन तद्विशेषणम्। विशेषणे ऽर्थे लिङ्गात्तृतीया भवति। जाटाभिस्तापसो दृष्टः। खड्गेन भटो दृष्टः।

\sutra{कर्तिरि च}

कर्तर्यर्थे लिङ्गात्तृतीया भवति। रुद्रेण त्रिपुरं दग्धम्। विष्णुना बद्धो बलिः।

\sutra{कालभावयोः सप्तमी}

काले भावे च भावो धात्वर्थः तस्मिन्नर्थे लिङ्गात्सप्तमी भवति। शरदि पुष्प्यन्ति सप्तच्छदाः। वर्षासु पुष्प्यन्ति कुटजानि। भावे यथा अग्निहोत्रे हूयमाने गतः। गोषु दुह्यमानासु आगतः।

\sutra{स्वामीश्वराधिपतिदायादसाक्षिप्रतिभूप्रसूतैः षष्ठी च}

स्वामिन् ईश्वर अधिपति दायाद साक्षिन् प्रतिभू प्रसूत इत्येतैर्योगे सति लिङ्गात् षष्ठी भवति सप्तमी च। धनस्य स्वामी धने स्वामी। धनस्य ईश्वरः धने ईश्वरः। धनस्य अधिपतिः धने अधिपतिः। धनस्य दायादः धने दायादः। धनस्य प्रतिभूः धने प्रतिभूः। गवां साक्षी गोषु साक्षी। गवां प्रसूतः गोषु प्रसूतः।

\sutra{निर्धारणे च}

जात्या गुणेन क्रियया संज्ञया वापृथक्कृतिः समूहादेकदेशस्य निर्धारणमिति स्मृतम्। येषां मध्यात्कस्यचिन्निर्धारणं पृथक्करणं तेषां षष्ठी भवति सप्तमी च। पुरुषाणां मध्ये क्षत्रियः शूरः पुरुषेषु वा। गवां कृष्णागौः संपन्नक्षीरा गोषु वा। अध्वगानां मध्ये धावन्तः शीघ्रगाः अध्वगेषु वा। ब्राह्मणानां मध्ये देवदत्तः पाठकः ब्रह्मणेषु वा।

\sutra{षष्ठी हेतुप्रयोगे}

हेतुशब्दप्रयोगे सति यो हेतुस्तस्य षष्ठी भवति। हेत्वर्थे इति तृतीया प्राप्ता। अन्नस्य हेतोर्वसति। विद्याया हेतोर्वसति।

\sutra{स्मृत्यर्थकर्मणि}

स्मृतिरर्थो योषां ते स्मृत्यर्थाः तेषां कर्म \textbf{स्मृत्यर्थकर्म} तस्मिन्कर्मणि षष्ठी भवति। पितुः स्मरति। मातुर्ध्यायति। प्रियायाश्चिन्तयति।

\sutra{करोतेः प्रतियत्ने}

करोते डुकृञ् करणे इत्यस्य यत्कर्म तस्मिन्कर्मणि षष्ठी भवति। प्रतियत्ने गम्यमाने सति। सतो गुणान्तरोत्पादनं प्रतियत्नः। एधोदकस्योपस्कुरुते। असिपत्त्रस्योपस्कुरुते।

\sutra{हिंसार्थानामज्वरे}

\textbf{हिंसार्थानां} धातूनां यत्कर्म तस्मिन्कर्मणि षष्ठी भवति \textbf{अज्वरे} ज्वर रोगे इत्यस्य कर्म वर्जयित्वा। चौरस्य रुजति रोगः। चौरस्य तुदति तोदः। चौरस्यामयति आमयः। \textbf{अज्वरे} इति किम्? चौरं ज्वरयति ज्वरः।

\sutra{कर्तृकर्मणोः कृति नित्यम्}

कृति कृत्प्रत्ययान्ते शब्दे प्रयुज्यमाने सति कर्तरि कर्मणि च नित्यं षष्ठी भवति। भवतः आसिका। भवतः शायिका। भवतः अग्रगामिका। पर्यायार्हणेषु चेति स्त्रियां भावे वुण्। युवजामना कान्ता। अत्र कर्तरि षष्ठी। कर्मणि यथा वज्रस्य भर्ता। भुवनस्य रक्षिता। यवानां लावकः। सक्तूनां पायकः।

\sutra{न निष्ठादिषु}

निष्ठादिषु कृत्प्रत्ययेषु प्रयुज्यमानेषु सत्सु कर्तृकर्मणोः षष्ठी न भवति। पूर्वेण प्राप्ता। देवदत्तेन जितम्। देवदत्तेन कटः कृतः। 

\sutra{षडो णो ने}

षष्शब्दस्य अन्त्यस्य हशषच्छेति डत्वे कृते णत्वं भवति ने नकारे परे। षण्णाम्। षण्णवतिः। षण्णगरी। ने इति किम् षड्भिः।

\sutra{मनोरनुस्वारो धुटि}

मकारनकारयोरनुस्वारो भवति धुटि परे। रंस्यते। हंसि। धुटीति किम्? गम्यते। हन्यते।

\sutra{वर्ग्ये वर्गान्तः}

मकारनकारयोर्वर्गान्तो भवति वर्ग्ये धुटि परे। शान्तः। शङ्किता। शमु उपशमे गत्यर्थाकर्मकेति क्तः। पञ्चमोपधेत्युपधादीर्घः। धुटीति किम्? चम्नोति हन्मि चमु अदने हन हिंसागत्योः वर्तमाना ति मि नु स्वादेर्गुणः।

\sutra{तवर्गश्चटवर्गयोगे चटवर्गौ}

तवर्गः चटवर्गाभ्यां योगे सति यथासङ्ख्यं चटवर्गावादेशौ भवतः राज्ञः भृज्जति। याच्ञा। भ्रस्जपाके भुयाचृयाच्ञायाम्। याचिविच्छिप्रच्छीति नङ्। टवर्गयोगे ईट्टे अड्डति षण्णाम्। ईडस्तुतौ वत्तमाना ते अन्वि अदादेर्लुक्। अघोषेष्वटामितिडस्यटः अद्डअभियोगे अन्वि अनेन तस्य टः दस्यडः।

\sutra{नामिकरपरःप्रत्ययविकारागमस्थःसिष्षं नुविसर्जनीयषान्तरोपि}

नामिनः ककारात् रेफाच्चपरः सकारः प्रत्ययस्थो विकारस्थ आगमगस्थश्च षत्वमापद्यते नुविसर्जनीयषान्तरोपि नुकारागमेन विसर्जनीयेन षकारेण व्यवहितोपीत्यर्थः। नामिनः अग्निषु वायुषु कात् वाक्षु त्वक्षु रेफात् गीर्षु धूर्षु इति प्रत्ययस्थः एषः इति विकारस्थः सर्वेषामित्यागमस्थः नुविसर्जनीयषान्तरोपि यथा संर्पीषि सर्पिषुः सर्पिष्षु। नामिन इति किम्? असौ।

\sutra{रषृवर्णेभ्यो नो णमनन्त्यः स्वरहयवकवर्गपवर्गान्तरोऽपि}

रेफषकारऋवर्णेभ्यः परो नकारः अनन्त्यः अपदान्तः णकारमापद्यते स्वरेण हकारेण यकारेण वकारेण कवर्गेण पवर्गेण च व्यवहितोऽपि इत्यर्थः। रेफात् चतुर्णाम्। षात् पूष्णाम्। ऋवर्णात् नृणां वा नॄणाम्। स्वरहयवकवर्गपवर्गैः केवलैः समुदितैश्च व्यवहितो। यथा हरेण पुरुषेण अहेण आर्येण गर्वेण अर्केण मूर्खेण सर्गेण अर्घेण सर्पेण रेफेण दर्भेण धर्मेण। रषृवर्गणेभ्य इति किम्? धनम्। अनन्त्य इति किम्? नरान्। वृक्षान्। पितॄन्।

\sutra{स्त्रियामादाप्}

आत् अकारान्तात् लिङ्गात् स्त्रियां वर्तमानात् आप्प्रत्ययो भवति। पकारः स्वरार्थः। देवदत्ता या सा का। स्त्रियामिति किम्? कीलालपाः। आदिति किम्? मतिः।

\sutra{नदाद्यन्च्वाह्वन्सन्तृसखिनान्तेभ्य ई}

नदादिभ्यः शब्देभ्यः तथा अन्च् वाह् उ अन्स् अन्त् ऋ सखि नान्त एभ्यः स्त्रियां वर्तमानेभ्यः ईप्रत्ययो भवति। नदादिभ्यस्तावत् नदी मही गौरी अन्च् प्राची प्रतीची तिरश्ची उदीची वाह् पष्टौही तुर्यौही। उ पट्वी मृद्वी। अन्स् पोचषी विदुषी। अन्त् पचन्ती देव्यन्ती। ऋ कर्त्री हर्त्त्री। सखि सखी। नान्त दण्डिनी पद्मिनी येभ्यः शब्देभ्यः स्त्रियामीप्रत्ययो दृश्यते तेन दादयः।

\sutra{ईकारे स्त्रीकृतेऽल्लोप्यः}

स्त्रियां कृते ईकारे परे अकारो लोपनीयो भवति। देवी गौरी मही। अनेन अकारलोपः। स्त्रीकृते इति किम्? कुण्डे।

\sutra{स्वरोह्रस्वोनपुंसके}

नपुंसकविषये योन्त्यः स्वरस्तस्य ह्रस्वो भवति सेनानि कुलम्। यवलुकुलम् नपुंसकेइति किम्? ग्रामणीः ना गौः।

इति लघुवृत्तौ नामप्रकरणे कारकपादश्चतुर्थः॥

\subsection{समासपादः}

\sutra{नाम्नां समासो युक्तार्थः}

नामसंज्ञां लिङ्गस्य विभक्त्यन्तस्य पूर्वाचार्याः कृतवन्तः।

अष्टौ यत्र प्रयुज्यन्ते नानार्थेषु विभक्तयः। \\
तन्नामकवयः प्राहुर्भेदे वचनलिङ्गयोः॥

\textbf{नाम्नां समासः} समसनं पृथक्स्थितानां एकीभावो भवति यदिपूर्वपदोत्तरपदानां \textbf{युक्तार्थः} परस्परसम्बन्धोऽर्थः भवति। नीलं च उत्पलं नीलोत्पलम्। राज्ञः पुरुषः राजपुरुषः। युक्तार्थ इति किम्? वस्त्रं नीलम् उत्पलं वर्षासु पुष्प्यति। भार्या राज्ञः पुरुषो देवदत्तस्य।

\sutra{तत्स्था लोप्या विभक्तयः}

तत्र समासाऽर्थे वाक्ये याः स्थिता विभक्तयः ताः लोप्या भवन्ति। राज्ञः पुरुषः राजपुरुषः। राजपुरुष इति षष्ठीतत्पुरुषः शङ्कुलया खण्डः शङ्कुलाखण्डः। यूपाय दारु यूपदारु। वृकेभ्यो भयं वृकभयं। लोप्या इति अर्हप्रत्ययनिर्देशादन्यस्यापि लोपार्हस्य यथादर्शनं लोपः। देवमुखमिव मुखं यस्य सः देवमुखः। घृतपूर्णो घटः घृतघटः। अविद्यमानं पापं यस्य सः अपापः। अत्र मुखपूर्णविद्यमानानां शब्दानां लोपः।

\sutra{प्रकृतिश्च स्वरान्तस्य}

स्वरान्तस्य शब्दस्य समासे विभक्तिलोपे च कृते प्रकृतिर्यथास्वरूपं भवति प्रत्ययलोपे प्रत्ययलक्षणम् इति वचनाद्यत्प्राप्तं तन्न भवति पितरि निपुणः पितृनिपुणः अत्र अर्ङाविति अर् प्राप्तः सखायं प्राप्तः सखिप्राप्तः। अत्र घुटि त्वै इति ऐत्वं प्राप्तम्।

\sutra{व्यञ्जनान्तस्य यत्सुभोः}

व्यञ्जनान्तस्य शब्दस्य सुप्रत्यये भकारे च परे यत्कार्यं प्राप्तं तत्समासे विभक्तिलोपे च कृते भवति। विदुषः अर्थः विद्वदर्थः। अत्र न लोपः दत्वं च भवति। राज्ञः अर्थः राजार्थः। अत्र न लोपः। दिवं गतः द्युगतः। अत्र उत्वम्। षण्णामाश्रयः षडाश्रयः अत्रडत्वम्।

\sutra{पदे तुल्याधिकरणे विज्ञेयः कर्मधारयः}

यत्र समासे \textbf{तुल्याधिकरणे} एकद्रव्यविषये उत्तर\textbf{पदे} परे पूर्वपदं समस्यते स समासः \textbf{कर्मधारय}संज्ञो विज्ञातव्यः। ब्राह्मणी चासौ दारिका ब्राह्मणदारिका। कर्मधारयसंज्ञे तु पुंवद्भावो विधीयते इति पूर्वपदस्य पुंवद्भावः। \textbf{तुल्याधिकरणे} इति किम्? कल्याण्या दुहिता कल्याणीदुहिता। कर्मधारयसंज्ञया क्व प्रयोजनम्? कर्मधारयसंज्ञे तु पुंवद्भावो विधीयते इत्यत्र।

\sutra{संख्यापूर्वो द्विगुरिति ज्ञेयः}

\textbf{सङ्ख्यापूर्वः} समासः \textbf{द्विगु}संज्ञो \textbf{ज्ञेयः} ज्ञातव्यः। तद्धितानां ये अर्थास्तेषु सत्सु उत्तरपदे परे समाहारे चाभिधेये द्विगुरिष्यते। तद्धितार्थेषु तावत् पञ्चसु कपालेषु संस्कृतः पुरोडाशः पञ्चकपालः। अत्र संस्कृतेऽर्थे एवमादेरणिष्यते इत्यण्। तस्य द्विगोरलुगनपत्य इति तन्त्रान्तरात् लुक्। उत्तरपदे पञ्च गावो धनं यस्य सः पञ्चगवधनः। अत्र धनशब्दे उत्तरपदे पञ्चगव इति द्विगुसमासान्तगतेति गोशब्दादकारस्समासान्तः। पञ्चानां गवां समाहारः पञ्चगवम्। पूर्ववदकारसमासान्तः। पञ्चानां नदीनां समाहारः पञ्चनदम्। संख्याया नदीगोदावर्योरिति नदीशब्दादकारः समासान्तः इवर्णावर्णाविति ईकारलोपः द्विगुसंज्ञया क्व प्रयोजनम्? द्विगोरित्यत्र।

\sutra{तत्पुरुषावुभौ}

उभौ कर्मधारयः द्विगुश्च तत्पुरुषौ भवतः। कुत्सितश्चासौ अश्वः कदश्वः। एवं कदर्यः कर्मधारयस्य तत्पुरुषत्वात् कोः कत् इति कत् आदेशः। पञ्चानां गवां समाहारः पञ्चगवम्। द्विगुसमासस्य तत्पुरुषत्वाद्गोरतद्धितलुकीति टच् समासान्तः। सखिराजशब्दाभ्यां तत्पुरुष एव स्मर्यते। पञ्चानां राज्ञां समाहारः पञ्चराजी। द्विगोरिति ईप्रत्ययः। ईकारे स्त्रीकृतेल्लोप्य इति अकारलोपः।

\sutra{विभक्तयो द्वितीयाद्या नाम्ना परपदेन तु समस्यन्ते समासो हि ज्ञेयस्तत्पुरुषस्स च}

यत्र समासे \textbf{द्वितीयाद्या विभक्तयो} द्वितीयादिविभक्तान्तानि पदानि \textbf{परपदेन नाम्ना समस्यन्ते स समासस्तत्पुरुषसंज्ञो ज्ञेयः} ज्ञातव्यः। धर्मं श्रितः धर्मश्रितः। शङ्कुलया खण्डः शङ्कुलाखण्डः। यूपाय दारु यूपदारु। वृकेभ्यो भयं वृकभयम्। राज्ञः पुरुषः राजपुरुषः। अक्षेषु शौण्डः अक्षशौण्डः। परपदेनेतिकिम्? गतो ग्रामम्। तत्पुरुषसंज्ञया क्व प्रयोजनम्? नस्य तत्पुरुषे लोप इत्यत्र।

\sutra{स्यातां यदि पदे द्वे तु यदि वा स्युर्बहून्यपि तान्यन्यस्य पदस्यार्थे बहुव्रीहिः}

यत्र समासे द्वे पदे स्यातां यदि वा बहूनि पदानि स्युः तानि पदानि अन्यस्य पदस्यार्थे वाच्ये सति समस्यन्ते स समासो बहुव्रीहिसंज्ञो भवति। दण्डैश्च दण्डैश्च यद्युद्धं तद्दण्डादण्डि इच्कर्मव्यतिहारे इति इच् समासान्तः। आकारो महतः कार्य इति कार्यग्रहणादाकारः इवर्णावर्णलोपः अव्ययं एतत् अव्ययाच्चेति सिलोपः। आरूढो वानरो यं वृक्षं सः आरूढवानरो वृक्षः। सहपुत्रेण वर्तते यः स सपुत्रः। सहभार्यया वर्तते यः स सभार्यः। ङिरौसपूर्व इति ज्ञापकात्सहशब्दस्य सभावः गोस्त्रियोरुपसर्जनस्येति उत्तरपदस्य ह्रस्वः। कृतः उपकारो येन स कृतोपकारः। रोचितो मोदको यस्मै स रोचितमोदकः। उद्धृतः ओदनो यस्या सा उद्धृतौदना स्थाली। चित्रा गावो यस्य स चित्रगुः। पुंवद्भाषितपुंंस्का इति चित्राशब्दस्य ह्रस्वः। शुद्धं जलं यस्याः सा शुद्धजला नदी। आसन्नाः दशानां ये ते आसन्नदशाः। एवं आसन्नविंशः आसन्नत्रिंशः। बहुव्रीहौ संख्येये डच् बहुगणादिति डच्। समासान्तः डानुबन्धेन्त्यस्वरादेर्लोपः बहूनि पदानि यथा। मत्ताः बहवो मातङ्गाः यत्र तत् मत्तबहुमातङ्गं वनम्। बहुव्रीहिसंज्ञया क्व प्रयोजनम्? बहुव्रीहौ चेत्यादिषु।

\sutra{विदिक्तथा}

विदिग्यत्र अन्यपदार्थाः स समासो बहुव्रीहिसंज्ञो भवति। दक्षिणस्याश्च पूर्वस्याश्च दिशो र्यदन्तरालं विदिक् सा दक्षिणपूर्वा। एवं उत्तरपूर्वा पूर्वोत्तरा पूर्वदक्षिणा दक्षिणपश्चिमा इत्याद्या उदाहार्याः। सर्वनाम्नो वृत्तिमात्रेण इति पूर्वपदानां पुंवद्भावः।

\sutra{द्वन्द्वः समुच्चयो नाम्नोर्बहूनां वापि यो भवेत्}

एकस्यां क्रियायां अनेनकर्मादिसम्बन्धः समुच्चयः। यत्र समासेद्वयोर्नाम्निोः बहूनां वा नाम्नां समुच्चयः स समासो द्वन्द्वसंज्ञो भवति। स च द्वयोः अर्थयोरिष्यते। समाहारे इतरेतरयोगे च। यत्र समूहप्राधान्यं स समाहारः। यत्र पृथक्स्थितानां पदार्थानां स इतरेतरयोगः। पदार्थानां नित्यसंयोगो वा तत्र प्रायशः। प्राण्यङ्गानां सेनाङ्गानां तूर्याङ्गानां येषामस्थिरहितं शरीरं दृश्यते तेषां क्षुद्रजन्तूनां च समूहेनैव स्थितिदर्शनात्समाहार एव द्वन्द्वो भवति। पाणी च पादौ च तत्पाणिपादम्। शिरोग्रीवम्। जङ्घोरु। इति प्राण्यङ्गानाम्। तूर्याङ्गानां यथा मार्दङ्गिपाणविकम् वांशिकवैणविकम्। सेनाङ्गानां यथा हस्त्यश्वम् रथिकाश्वारोहम्। क्षुद्रजन्तूनां यथा यूकालिक्षम् दंशमषकम् अप्राणिनामपि फलादीनाम् प्रायशः समूहवृत्तित्वात्समाहार एव। बदरामलकम् आराशस्ति स्वरो ह्रस्वो नपुंसके इति ह्रस्वः। प्राणिनां तु पृथग्वृत्तित्वादितिरेतरयोग एव ब्राह्मणक्षत्रियौ। नित्यवैरिणां तुपरस्परविरोधात्समूहवृत्तित्वाभावेऽपि समाहार एव। येषां च विरोधः शाश्वतिक इति तन्त्रान्तरात् अहिनकुलम् मूषिकमार्जारम् बिडालोन्दुरुम्। एवमन्यत्रापि लक्ष्यानुसारेण समाहारे इतरेतरयोगे द्वन्द्वसमासो ज्ञेयः।

\sutra{अल्पस्वरतरं तत्र पूर्वम्}

तत्र द्वन्द्वे यदल्पस्वरतरं तत्पूर्वमेव प्रयोज्यम्। धवखदिरौ प्लक्षन्यग्रोधौ तण्डुलकिण्वे उलूखलमुसले इत्यादौ पूर्वनिपातः। अर्चितत्वात्तरग्रहणाल्लघ्वक्षरस्य पूर्वानिपातः। कुशकाशम् अग्निसज्ञकस्य पूर्वनिपातः द्वन्द्वे घि इति तन्त्रान्तरात् अग्नीषोमौ ईदग्नेः सोमवर्णयोरितितन्त्रान्तरात् ईत्वम् अग्नेस्तुत्स्तोमसोम इति षत्वम्। 

\sutra{यच्चार्चितं द्वयोः}

द्वयोः नाम्नोः द्वन्द्वे यदर्शितं पूजितं तत्पूर्वमेव प्रयोज्यम् वासुदेवार्जुनौ मातापितरौ पितापुत्रौ। आनङ्यतो द्वन्द्व इति पूर्वपदान्त्य ऋकारस्यात्वम् ब्राह्मणक्षत्रियौ सूर्याचन्द्रमसौ देवताद्वन्द्वे चेति पूर्वपदान्त्यस्यात्वम्। जायापती भार्यापती अत्रापि पतिर्जायां प्रविशति गर्भो भूत्वेह जायते जायायाः तद्विजायात्वं यदस्यां जायते पुनरिति स्मरणाज्जायायाः अर्चितत्वम्। चशब्दात् क्वचिदनर्चितस्यापि पूर्वनिपातः। चन्द्रार्कौ शुक्रबृहस्पती कालकामौ काकमयूरौ शत्रुमित्रे इत्यादि द्वयोरितिकिं? सर्वज्वरान् घ्नन्तु तवानिरुद्रप्रद्युम्नसङ्कर्षणवासुदेवाः इति।

\sutra{पूर्वं वाच्यं भवेद्यस्य सोव्यययीभाव इष्यते}

यस्य समासस्य पूर्वं पूर्वपदार्थो वाच्यं प्रधानं भवेत् स समासः अव्ययीभावसंज्ञ इष्यते। समीपं कुम्भस्य उपकुम्भम्। समृद्धिः कर्तृणां सुकर्तृ। अभावः पापानाम् अपापम्। अत्ययः शीतानाम् अतिशीतम्। पश्चाद्रथस्य अनुरथम्। विधिमनतिक्रम्य यथाविधि। स्थाने स्थाने यथास्थानम्। योयःस्वः यथास्वम्। येन येन प्रकारेण यथायथम्। इतिनिपातनात् आकारस्य ह्रस्वः। अवययीभावसंज्ञया क्व प्रयोजनम् अव्ययीभावादकारान्ताद्विभक्तीनामपञ्चम्या इत्यादिषु।

\sutra{स नपुंसकलिङ्गः स्यात्}

सोव्ययीभावसमासो न पुंसकलिङ्गो भवति। उक्तान्युदाहरणानि उपकुम्भमित्यादीनि।

\sutra{द्वन्द्वैकत्वम्}

द्वन्द्वस्य यदेकत्वं समाहारस्तत् नपुंसकलिङ्गं भवति पाणिपादम् बदरामलकम्।

\sutra{तथा द्विगोः}

द्विगोः समासस्य नपुंसकत्वं भवति तथानेनैव प्रकारेण येनैवद्वन्द्वस्येत्यर्थः। पञ्चानां नदीनां समाहारः पञ्चनदम्। दशानां गवां समारारः दशगवम्। सङ्ख्याया नदीगोदावर्योः गोरतद्धितलुकीति च। अकारः समासान्तः। समाहारादन्यत्र यो द्विगुस्तस्य न भवति। पञ्चसु कपालेषु संकृतः पुरोडाशः पञ्चकपालः।

\sutra{पुंवद्भाषितपुंस्कारूनूङपूरण्यादिषु स्त्रियां तुल्याधिकरणे}

भाषितपुंस्कायास्त्रीविशेषणं स्त्रीलिङ्गं अनूङ् रूङ्प्रत्ययरहितं मास्त्रीलिङ्गे पूरण्यदिवर्जिते तुल्याधिकरणे टाकद्रव्यविषये तुत्तरपदेपरेपुंवद्भवति पुंसियद्रूपंतदापद्यते इत्यर्थः दर्शनीयाभार्गायस्यसदर्शनीयभार्यः उत्तरपदस्यगोस्त्रियोरुपसर्जनस्य इति ब्रस्वः कल्याणीभार्याययम्यसकल्याणभार्यः एवं स्वरूपभार्यः भाषितपुंस्केतिकिम्? खट्वाभार्यायस्यसखट्वाभार्यः स्त्रीलिङ्गनिर्देशः किम्? अतिरिक्तलंगतिर्यस्यस अतिरिगतिः अतिनुजलं गतिर्यस्य स अतिनुगतिः पुंवद्भावेन ह्रस्वनिवृत्तिः स्यात् अनूज् इति किम्? ब्रह्मबन्धूभार्या यस्य स ब्रह्मबन्धूभार्यः ब्रह्मबन्धुशब्दात् रूङुतइत्प्रङ् ङकारः इहार्थः अपूरण्यदिष्वितिकिम्? कल्याणीपङ्चमीयासांताः कल्याणीपञ्चमारात्रयः अम्पूरण्योप्रामाण्योरित्यप् समासान्तः इवर्णावर्णलोपः स्त्रियामिति किम्? कल्याणीप्रधानं यासां ताः कल्याणीप्रधानाः तुल्याधिकरणे इति किं? कल्याण्यदुहिताकल्याणीदुहिताः।

\sutra{संज्ञापूरणीकोपधात्तु न}

संज्ञाभूतापूरणप्रत्ययान्ताककारोपधा च भाषितपुंस्का पुंवन्न भवति स्त्रीलिङ्गे तुल्याधिकरणे पदे पूर्वेणप्राप्तिः दत्ताभार्यायस्य स दत्ताभार्यः पाचिकाभार्यः पञ्चमीभार्यः नागरिकाभार्यः गाषिकाभार्यः नर्तकीभार्यः।

\sutra{कर्मधारयसंज्ञे तु पुंवद्भावो विधीयते}

संज्ञापूरणीकोपधानां कर्मधारयसंज्ञे समासे तुल्याधिकरणे स्त्रीलिङ्गे पूरण्यादिवर्जिते उत्तरपदेपरे पुंवद्भावो विधीयते पूर्वेण निषेधः यः प्राप्तः दत्ताचासौवृन्दारिका दत्तवृन्दारिका। पञ्चमवृन्दारिका।

\sutra{आकारोमहतःकार्यस्तुल्याधिकरणे पदे}

महदित्यस्य आकारः कार्यः तुल्याधिकरणे पदे परे महांश्चासौ देवः महादेवः। महान् कायो यस्य स महाकायः तुल्याधिकरणे इति किम्? महतः पुत्रो महत्पुत्रः। कार्यग्रहणादन्यत्रापि पूर्वपदस्यात्वं भवति दण्डैश्च दण्डैश्च यद्युद्धं तद्दण्डादण्डि मातापितरौ पितापुत्रौ एकादश।

\sutra{नस्य तत्पुरुषे लोपः}

नस्यपदस्य तत्पुरुषे प्रसिद्धस्य उत्तरपदे परे नकारस्य लोपो भवति नब्राह्मणः अब्राह्मणः नवृषलः अवृषलः पदस्येति किम्? वामनपुत्रः।

\sutra{स्वरेक्षरविपर्ययः}

अतरयोर्विपर्ययः अक्षरविपर्ययः नस्यपदस्य अक्षरविपर्ययीभवति स्वरे परे न अश्वः अनश्वः न उष्ट्रः अनुष्ट्रः।

\sutra{कोः कत्}

कुशब्दस्य तत्पुरुषे प्रसिद्धस्य कत् आदेशो भवति स्वरे परे कुत्सितश्चासौवश्च कदश्वः तत्पुरुषे इति किम्? कुत्सित ऊष्मा अण्डे यस्य स कूण्माण्डः।

\sutra{का त्वीषदर्थे}

ईषदर्थे वर्तमानस्य कुशब्दस्य का आदेशो भवति ईषन्मधुरम् कामधुरम् एवं कालवणं काम्बम् ईष्दर्थे इति किम्? कुत्सितो देशः कुदेशः।

\sutra{अक्षे}

अक्षशब्दे उत्तरपदे परे कुशब्दस्य काआदेशो भवति। अनीषदर्थं वचनम् कुत्सितः अक्षः काक्षः।

\sutra{पुरुषे तु विभाषया}

कुशब्दस्य पुरुषशब्दे परे काआदेशो भवति विभाषया विकल्पेन कुत्सितः पुरुषः कापुरुषः कुपुरुषो वा।

\sutra{याकारौस्त्रीकृतौह्रस्वौक्वचित्}

स्त्रियां कृतौ स्त्रीकृतौ ईकारआकारौ उत्तरपदे परे ह्रस्वौ भवतः क्वचित् न सर्वत्र रेवति मित्रः रोहिणिदत्तः इष्टकाइषीका मालाएषांयथासंख्यं सचितसूलभारिन् एषु उत्रपादेषु परेषु अन्त्यस्य ह्रस्वो भवति इष्टकेषीकामालानां चितरुलभारिथितितन्त्रान्तरात् इष्टकामिश्चितम् इष्टकाचितम् इषीकायास्थूलं इषीकरूलम् मालां बिभर्तीति मालभारिणीकत्या क्वचिदिति वचनात् इहन भवति लक्ष्मीतरातन्त्रीकल्पा।

\sutra{ह्रस्वस्य दीर्घता}

ह्रस्वस्य क्वचिद्दीर्घता भवति दीर्घो भवतीत्यर्थः अष्टसुकपालेषु संस्कृतः पुरोडाशः अष्टाकपालः द्व्यङ्गुलौकर्णौ यस्य स द्व्यङ्गुलाकर्णः एवं द्विगुणाकर्णः क्वचिदितिवचनादिहनभवति छिन्नकर्णः भिन्नकर्णः स्रुवकर्णः स्वस्तिकर्णः।

\sutra{अनव्ययविसृष्टस्तुसकारं कपवर्गयोः}

अनव्ययशब्दस्य विसृष्टो विसर्जनीयः कपवर्गयोः परयोः क्वचित्सकारमापद्यते जिह्वामूलीयोपध्मानीययोर्वाधः अनुन्तरपदस्थादकारात्परस्य विसर्जनीयस्य कृकमिकं स कुम्भपात्रकुशाकर्ण एषु परेषु समासे सतिः अतः कृकमिकं स कुम्भाकुशाकर्णेष्वनव्ययस्येति तन्त्रान्तरात् कृ अयस्कारः कामि अयस्कामः कंस अयस्कंसः कुम्भ अयस्कुम्भ पात्र अयस्पात्रम् कुशा अयस्कुशा कर्ण अयस्कर्णः उत्तरपदस्थात्तु अकारात् परस्य विसर्जनीयस्य न भवति उपयःकामः आकारादपि परस्य क्वचित्स्मर्यते भास्करः कपवर्गयोरिति कखपफा एव गृह्यन्ते तेनेह न भवति अयोघनः पयोबिन्दुः अनव्ययविसृष्ट इति किम्? स्वःकामः। अन्तःकामः।

इति लघुवृत्तौ नामप्रकरणे समासपादः पञ्चमः॥

\subsection{तद्धितपादः}

\sutra{वाणपत्ये}

अपत्येऽर्थे लिङ्गादण् प्रत्ययो वा भवति। उपगोरपत्यम् औपगवः उपगुशब्दादण् णकारो वृद्धिरादौ सणे इत्यर्थः। आदिवृद्धिः उवर्णस्त्वोत्वमापाद्य इति ओत्वम्। पक्षे उपगोरपत्यं उपग्वपत्यमिति समासवाक्ये भवतः। संख्यासंभद्रपूर्वायामातुरणि परे ऋकारस्य उरादेशो भवति। मातुरुत्संख्यासंभद्रपूर्वाया इति तन्त्रान्तरात् षाण्मातुरः सांमातुरः भाद्रमातुरः। कन्याया अणि परे कनीनादेशो कानीनो व्यासः कानीनः कर्णः।

\sutra{ण्य गर्गादेः}

गर्गादेर्गणादपत्येऽर्थे ण्यप्रत्ययो वा भवति। गर्गस्यापत्यं गार्ग्यः वात्स्यः। पक्षे गर्गस्यापत्यं गर्गापत्यमिति। पौत्रादावपत्ये प्रत्ययोऽयं स्मर्यते तेन गर्गस्य पुत्रो गार्गिः इण्णतः इति इण्।

\sutra{पत्यन्ताश्च}

प्रजापतेरपत्यं प्राजापत्यः गार्हस्पत्यः दित्यदित्योरेयण्च दैत्यः दैतेयः आदित्यः आदितेयः गर्गवत्सवजसङ्कृतिव्याघ्रपात् शाण्डिल्यमुद्गलविदभृत् इत्याद्या गर्गादयः।

\sutra{कुञ्जादेरायनण्स्मृतः}

कुञ्जादेर्गणात् अपत्येऽर्थे आयनण् प्रत्ययोभवति। तदन्ताच्च स्वार्थे ण्यप्रत्यय इष्यते। स च स्त्रियां बहुत्वेनेष्यते। कौञ्जायन्यः कौञ्जायन्यौ कौञ्जायनाः कौञ्जायनी। एवं ब्राध्नायन्यः ब्राध्नायन्यौ ब्राध्नायनाः ब्राध्नायनी जातेरयोपधादिति ईप्रत्ययः कुञ्जब्रध्नरणनभचरलिगुलोह इत्याद्या कुञ्जादयः।

\sutra{स्त्र्यत्र्यादेरेयण्}

स्त्रीप्रत्ययान्तेभ्यः अत्र्यादिभ्यश्च अपत्येर्थे एयण्प्रत्ययो भवति सुपर्णाया अपत्यं सौपर्णेयः वैनतेयः गाङ्गेयः नद्या अपत्यं नादेयः एवं माहेयः रौहिणेयः अत्र्यादिभ्यः आत्त्रेयः नैधेयः पितृस्वसृमातृस्वसृभ्यामेयण् ईवर्णौ इष्येते एयणि वानुलोपश्च पितृष्वसुरपत्यं पैतृश्वस्रेयः पैतृश्वसेयः पैतृश्वस्रीयः। एवं मातृश्वस्रेयः मातृश्वसेयः मातृश्वस्रीयः  कल्याण्यादीनामेयणिपरे अन्त्यस्य इनादेशो भवति कल्याण्यादीनामिनङ्ङितितन्त्रान्तरात् कल्याण्या अपत्यं काल्याणिनेयः सुभगाया अपत्यं सौभागिनेयः हृद्भगसिन्ध्वन्ते पूर्वपदस्येति तन्त्रान्तादुभयपदवृद्धिः अत्त्रिनिधिशुद्धविष्टिमृकाण्डकृक्रतास इत्याद्या अत्त्र्यादयः।

\sutra{इण्णतः}

अतः अकारान्तात् लिङ्गादपत्येर्थे इण् प्रत्ययो भवति दक्षस्यापत्यं दाक्षिः प्लाक्षिः।

\sutra{बाह्वादेश्चविधीयते}

बाह्वादेर्गणात् अपत्येर्थे इण् भवति बाहोरपत्यं बाहविः औपचाकविः औपबिन्दविः बालाकिः सौमित्त्रिः औडुलोमिः शारलोमिः पौष्करसादिः उडुलोमन् शरलोमन् इत्येतयोः नस्य तुक्वचिदिति नलोपः पुष्करसद् शब्दस्य अनुशतिकादीनां चेत्युभयपदवृद्धिः जाङ्घिः दैवशर्मिः सुधातृव्यासवरुटनिषादचण्डालबिम्बानामिनिपरे अन्त्यस्य अक आदेश इष्यते सुधातुरपत्यं सौधातकिः वैयासकिः य्वौर्वृद्धिरागम इति यकारात्पूर्वम् ऐकारागमः वारुटकिः नैषादकिः चाण्डालकिः बैम्बकिः बाहु उपचाक्वः उपबिन्दुः बलाका सुमित्त्रा उडुलोमन् शरलोमन् पुष्करसद् जङ्घादेवशर्मन् इत्यादयोबाह्वादयः।

\sutra{रागान्नक्षत्रयोगाच्च समूहात्सास्य देवता।\\ तद्वेत्त्यधीते तस्येदमेवमादेरणिष्यते}

रागात् राज्यतेयेन स रागः रागवाचिनः शब्दात् अण् भवति। तेन रक्तमित्यत्रार्थे कषायेण रक्तं काषायम् कुसुम्भेन रक्तं कौसुम्भम् कौङ्कुमम् हारितालम्।

नक्षत्रयोगाच्च नक्षत्रेण योगेमम्यमाने नक्षत्रयोगात् अण् भवति पुष्येण युक्ता पौर्णमासी पौषी तैषी तिष्यपुष्ययोर्नक्षत्राणीति तन्त्रान्तरात् यलोपः सापौर्णमासी यस्य मासस्य स पौषः तैषः सास्मिन्पौर्णमासीति संज्ञाया मिति ईप्रत्ययान्तादण् इवर्णावर्णलोपः एवं मघाभिर्युक्ता पौर्णमासी फाल्गुणीमा यस्य मासस्य स फाल्गुणः चित्रयायुक्तापौर्णमासी चैत्री सा यस्य मासस्य स चैत्रः विशाखया युक्ता पौर्णमासी वैशाखी सा यस्य मासस्य स वैशाखः ज्येषठया युक्ता पौर्णमासी ज्यैष्ठी सा यस्य मासस्य सज्यैष्ठः आषाढाभिर्युक्ता पौर्णमासी आषाढी सा यस्य मासस्य स आषाढः श्रवर्णेन युक्ता पौर्णमासी श्राणी सायस्य मासस्य स श्रावणः भद्रपदाभिर्युक्तापौर्णमासी भाद्रपदी सायस्य मासस्य स भाद्रपदः अश्वयुग्भ्यां युक्तापौर्णामासी आश्वयुजी सा यस्य मासस्य स आश्वयुजः कृत्तिकाभिर्युक्ता पौर्णमासी कार्तिकी सायस्य मासस्य स कार्तिकः मृगशीर्षेणयुक्ता पौर्णमासी मार्गशीर्षी सा यस्य मासस्य स मार्गशीर्षः।

समूहात् समूहेर्थे लिङ्गात् अण् भवति कपोतानां समूहः कापोतम्। शुकानां समूहः शौकम्। मायूरम्। तित्तिरीणां समूहः तैत्तिरम् गर्भिणीनां समूहः गार्भिणम् इवर्णावर्णलोपः।

सास्य देवता लिङ्गात् सा अस्य देवता इत्यत्रार्थे अण् भवति शिवोदवता अस्येति शैवः विष्णुर्देवतास्येति वैष्णवः एवं बौद्धः। 

तद्वेत्त्यधीते लिङ्गात् तद्वेत्तितदधीते इत्यर्त्रार्थे अण् भवति छन्दोवेत्त्यधीतेवा छान्दसः व्याकरणं वेत्त्यधीते वा वैयाकरणः य्वोर्वृद्धिरागमइति यकारात्पूर्वऐकारागमः।

तस्येदम् लिङ्गात् तस्येदमित्यत्रार्थे अण् भवति मृगस्येदं मार्गं चर्म वेणोरयं वैणवो दण्डः।

एवमादेरणिष्यते एवमादावित्यत्रार्थे अण् भवति तत्र भवः स्रुघ्ने भवः स्रौघ्नः एवं कान्यकुब्जः मथुरायां भवः माथुरः तत आगतः स्रुघ्नादागतः स्रौघ्नः तत्र जातः स्रघ्ने जातः स्रौघ्नः।

\sutra{तेन दीव्यति संसृष्टं तरतीकण् चरत्यपि।\\ पण्याच्छिल्पान्नियोगच्च क्रीतादेरायुधादपि}

तेन दीव्यति लिङ्गात्तेन दीव्यतीत्यत्रार्थे इकण् भवति तेन दीव्यति अक्षैर्दीव्यति आक्षिकः शलाकाभिर्दीव्यति शालाकिकः।

संसृष्टम् तेन संसृष्टमित्यत्रार्थे इकाण् भवति दध्ना संसृष्टं दाधिकम्।

तरति लिङ्गात्तरतीइत्यत्रार्थे इकण् भवति नावा तरति नाविकः उडुपेन तरति औडुपिकः वैणुभारिकः।

इकण्चरत्यपि लिङ्गात्तेन चरतीत्यत्रार्थे इकण् भवति। शकटेन चरति शाकटिकः आश्विकः औष्ट्रिकः चरतिर्जीवनेपि वर्तते वेतनेन चरति वैतनिकः व्यावहारिकः व्यापारिकः य्वोर्वृद्धिरागमोनेष्यते स्वागतादीनां चेतितन्त्रान्तरात्।

पण्यात् पण्यं विक्रेयमुच्यते पण्यवाचिनोलिङ्गात् तदस्य पण्यमित्यत्रार्थे इकण्भवति लवणं पण्यमस्य लावणिकः गन्धः पण्यमस्य गान्धिकः।

शिल्पात् लिङ्गात्तदस्यशिल्पमित्यत्रार्थे इकण् भवति मृदङ्गः शिल्पमस्य मार्दङ्गिकः पाणविकः।

नियोगाच्च नियुज्यते यत्र स नियोगः नियोगवाचिनः शब्दात् तत्रनियुक्त इत्यत्रार्थे इकण् भवति द्वारे नियुक्तः दौवारिकः द्वारादीनां चेति वकारात्पूर्वऔकारागमः।

क्रीतादेः क्रीदादावर्थे लिङ्गात् इकण् भवति तेन क्रीतं षष्ट्या क्रीतं षाष्टिकम् सपुत्याक्रीतं साप्ततिकम्।

आयुधात् तदस्यायुधमित्यत्रार्थे इकण्भवति खड्गा आयुधं यस्य स खाड्गिकः प्रासिकः।

\sutra{नावस्तार्येविषाद्वध्येतुलयासस्मितेपि च तत्र साधौ च यः}

नौशब्दात्तेन तार्यमित्यत्रार्थे विषशब्दात् तेन वध्य इत्यत्रार्थे तुलाशब्दात्तृतीयान्तात् सम्मिते चार्थे लिङ्गात्तत्र साधुरित्यत्रार्थे च य प्रत्ययो भवति नावा तार्यं नाव्यम् विषेणवश्यः विष्यः तुलयासम्मितं तुल्यम् अपिशब्दात् वयस् शब्दात् तेन तुल्य इत्यत्रार्थे च य प्रत्ययो भवति वयसातुल्योवयस्यः कर्मणिसाधुः कर्मण्यः शरणेसाधुः शरण्यः।

६२, १०५

\end{document}